\documentclass[11pt]{article}
% =======================================================================
%   Shared style for CS 300 notes.
%   Usage: place `% =======================================================================
%   Shared style for CS 300 notes.
%   Usage: place `% =======================================================================
%   Shared style for CS 300 notes.
%   Usage: place `\input{../../notes-style.tex}` right after \documentclass
%   in each chapter's lectures.tex and notes.tex.
% =======================================================================

% ---- Layout -----------------------------------------------------------
\usepackage[margin=1in]{geometry}
\usepackage{parskip}        % space between paragraphs, no indent
\usepackage{microtype}      % subtle kerning/spacing improvements

% ---- Colors (muted palette, easy on light-sensitive eyes) -------------
\usepackage{xcolor}
\definecolor{pagebg}       {HTML}{FAF6F0}  % warm cream background
\definecolor{bodyfg}       {HTML}{3D3A36}  % warm dark gray body text
\definecolor{heading}      {HTML}{3D6A6B}  % muted teal
\definecolor{subheading}   {HTML}{5B7A9C}  % dusty blue
\definecolor{accent}       {HTML}{8A6B4A}  % warm tan
\definecolor{linkcol}      {HTML}{5B7A9C}  % same as subheading
\definecolor{mutedtext}    {HTML}{6B6560}  % secondary / caption text
\definecolor{rulecol}      {HTML}{C8BFB0}  % separator / rule

% Callout box palette
\definecolor{defnFrame}    {HTML}{9DB89B}  \definecolor{defnBack}    {HTML}{EEF2E8}
\definecolor{ideaFrame}    {HTML}{9FB4CD}  \definecolor{ideaBack}    {HTML}{E8EEF5}
\definecolor{warnFrame}    {HTML}{CFAE87}  \definecolor{warnBack}    {HTML}{F5EDDF}
\definecolor{noteFrame}    {HTML}{BFA6AE}  \definecolor{noteBack}    {HTML}{F2E8EE}
\definecolor{exFrame}      {HTML}{9FBDBD}  \definecolor{exBack}      {HTML}{E8F0F0}
\definecolor{codeBack}     {HTML}{F3EDE2}
\definecolor{codeBorder}   {HTML}{D4CBBA}
\definecolor{codeKeyword}  {HTML}{6B4F7A}
\definecolor{codeString}   {HTML}{8A6B4A}
\definecolor{codeComment}  {HTML}{8A857F}

% ---- Page background + base text color --------------------------------
\usepackage{pagecolor}
\pagecolor{pagebg}
\color{bodyfg}

% ---- Fonts ------------------------------------------------------------
\usepackage[T1]{fontenc}
\IfFileExists{libertine.sty}{%
    \usepackage{libertine}      % warm, readable serif
    \usepackage[scaled=0.95]{inconsolata}  % softer monospace
}{}

% ---- Hyperref ---------------------------------------------------------
\usepackage{hyperref}
\hypersetup{
    colorlinks=true,
    linkcolor=linkcol,
    urlcolor=linkcol,
    citecolor=linkcol,
    pdfborder={0 0 0},
}

% ---- Section styling --------------------------------------------------
\usepackage[explicit]{titlesec}
\titleformat{\section}
    {\color{heading}\Large\bfseries}
    {\color{heading}\thesection\hspace{0.6em}}{0pt}{#1}
\titleformat{\subsection}
    {\color{subheading}\large\bfseries}
    {\color{subheading}\thesubsection\hspace{0.5em}}{0pt}{#1}
\titleformat{\subsubsection}
    {\color{accent}\normalsize\bfseries}
    {\color{accent}\thesubsubsection\hspace{0.5em}}{0pt}{#1}
\titleformat{\paragraph}[runin]
    {\color{accent}\bfseries}{}{0pt}{#1.}

\titlespacing*{\section}{0pt}{1.6em}{0.6em}
\titlespacing*{\subsection}{0pt}{1.2em}{0.4em}

% ---- Enumerate / itemize spacing --------------------------------------
\usepackage{enumitem}
\setlist[itemize]{leftmargin=1.4em, itemsep=2pt, topsep=4pt}
\setlist[enumerate]{leftmargin=1.6em, itemsep=2pt, topsep=4pt}

% ---- Code listings ----------------------------------------------------
\usepackage{listings}
\lstdefinestyle{notescode}{
    backgroundcolor=\color{codeBack},
    basicstyle=\ttfamily\small\color{bodyfg},
    keywordstyle=\color{codeKeyword}\bfseries,
    stringstyle=\color{codeString},
    commentstyle=\color{codeComment}\itshape,
    frame=single,
    framesep=6pt,
    rulecolor=\color{codeBorder},
    showstringspaces=false,
    breaklines=true,
    xleftmargin=10pt,
    xrightmargin=10pt,
    aboveskip=8pt,
    belowskip=8pt,
}
\lstset{style=notescode}

% ---- Callout boxes (tcolorbox) ----------------------------------------
\usepackage[most]{tcolorbox}

\newtcolorbox{defnbox}[1][]{
    enhanced, breakable,
    colback=defnBack, colframe=defnFrame, coltitle=bodyfg,
    title={\bfseries Definition \textemdash\ #1},
    fonttitle=\bfseries,
    boxrule=0.5pt, arc=2pt, left=8pt, right=8pt, top=6pt, bottom=6pt,
    attach boxed title to top left={xshift=10pt, yshift=-6pt},
    boxed title style={colback=defnFrame, colframe=defnFrame, arc=2pt},
    before skip=8pt, after skip=8pt,
}
\newtcolorbox{ideabox}[1][]{
    enhanced, breakable,
    colback=ideaBack, colframe=ideaFrame, coltitle=bodyfg,
    title={\bfseries Key idea #1},
    boxrule=0.5pt, arc=2pt, left=8pt, right=8pt, top=6pt, bottom=6pt,
    attach boxed title to top left={xshift=10pt, yshift=-6pt},
    boxed title style={colback=ideaFrame, colframe=ideaFrame, arc=2pt},
    before skip=8pt, after skip=8pt,
}
\newtcolorbox{warnbox}[1][]{
    enhanced, breakable,
    colback=warnBack, colframe=warnFrame, coltitle=bodyfg,
    title={\bfseries Gotcha #1},
    boxrule=0.5pt, arc=2pt, left=8pt, right=8pt, top=6pt, bottom=6pt,
    attach boxed title to top left={xshift=10pt, yshift=-6pt},
    boxed title style={colback=warnFrame, colframe=warnFrame, arc=2pt},
    before skip=8pt, after skip=8pt,
}
\newtcolorbox{notebox}[1][]{
    enhanced, breakable,
    colback=noteBack, colframe=noteFrame, coltitle=bodyfg,
    title={\bfseries Aside #1},
    boxrule=0.5pt, arc=2pt, left=8pt, right=8pt, top=6pt, bottom=6pt,
    attach boxed title to top left={xshift=10pt, yshift=-6pt},
    boxed title style={colback=noteFrame, colframe=noteFrame, arc=2pt},
    before skip=8pt, after skip=8pt,
}
\newtcolorbox{examplebox}[1][]{
    enhanced, breakable,
    colback=exBack, colframe=exFrame, coltitle=bodyfg,
    title={\bfseries Example #1},
    boxrule=0.5pt, arc=2pt, left=8pt, right=8pt, top=6pt, bottom=6pt,
    attach boxed title to top left={xshift=10pt, yshift=-6pt},
    boxed title style={colback=exFrame, colframe=exFrame, arc=2pt},
    before skip=8pt, after skip=8pt,
}

% ---- TikZ graph styles (added 2026-04-25 for ch_10 augmentation) ------
% tcolorbox[most] already loads tikz; this block declares the libraries
% + shared `vertex` / `edge` styles used by ch_10's general directed-graph
% diagrams. ch_9's tree-style TikZ uses inline overrides and is
% unaffected. Simple circle nodes + arrow edges only -- no production
% styling.
\usetikzlibrary{arrows.meta, positioning, calc}
\tikzset{
    vertex/.style={circle, draw, minimum size=7mm, inner sep=0pt,
                   font=\small, thick},
    vsource/.style={vertex, fill=ideaBack},
    vdone/.style={vertex, fill=defnBack},
    edge/.style={->, >=Stealth, thick},
    uedge/.style={-, thick},
    treeedge/.style={->, >=Stealth, thick},
    backedge/.style={->, >=Stealth, thick, dashed, red!70!black},
    fwdedge/.style={->, >=Stealth, thick, blue!60!black},
    crossedge/.style={->, >=Stealth, thick, green!50!black, dotted},
    mstedge/.style={-, very thick, accent},
}

% ---- Title block ------------------------------------------------------
\usepackage{titling}
\pretitle{\begin{center}\color{heading}\LARGE\bfseries}
\posttitle{\end{center}\vskip -0.4em}
\preauthor{\begin{center}\color{mutedtext}\normalsize}
\postauthor{\end{center}}
\predate{\begin{center}\color{mutedtext}\small}
\postdate{\end{center}\vspace{0.4em}{\color{rulecol}\hrule height 0.4pt}\vspace{1.2em}}
` right after \documentclass
%   in each chapter's lectures.tex and notes.tex.
% =======================================================================

% ---- Layout -----------------------------------------------------------
\usepackage[margin=1in]{geometry}
\usepackage{parskip}        % space between paragraphs, no indent
\usepackage{microtype}      % subtle kerning/spacing improvements

% ---- Colors (muted palette, easy on light-sensitive eyes) -------------
\usepackage{xcolor}
\definecolor{pagebg}       {HTML}{FAF6F0}  % warm cream background
\definecolor{bodyfg}       {HTML}{3D3A36}  % warm dark gray body text
\definecolor{heading}      {HTML}{3D6A6B}  % muted teal
\definecolor{subheading}   {HTML}{5B7A9C}  % dusty blue
\definecolor{accent}       {HTML}{8A6B4A}  % warm tan
\definecolor{linkcol}      {HTML}{5B7A9C}  % same as subheading
\definecolor{mutedtext}    {HTML}{6B6560}  % secondary / caption text
\definecolor{rulecol}      {HTML}{C8BFB0}  % separator / rule

% Callout box palette
\definecolor{defnFrame}    {HTML}{9DB89B}  \definecolor{defnBack}    {HTML}{EEF2E8}
\definecolor{ideaFrame}    {HTML}{9FB4CD}  \definecolor{ideaBack}    {HTML}{E8EEF5}
\definecolor{warnFrame}    {HTML}{CFAE87}  \definecolor{warnBack}    {HTML}{F5EDDF}
\definecolor{noteFrame}    {HTML}{BFA6AE}  \definecolor{noteBack}    {HTML}{F2E8EE}
\definecolor{exFrame}      {HTML}{9FBDBD}  \definecolor{exBack}      {HTML}{E8F0F0}
\definecolor{codeBack}     {HTML}{F3EDE2}
\definecolor{codeBorder}   {HTML}{D4CBBA}
\definecolor{codeKeyword}  {HTML}{6B4F7A}
\definecolor{codeString}   {HTML}{8A6B4A}
\definecolor{codeComment}  {HTML}{8A857F}

% ---- Page background + base text color --------------------------------
\usepackage{pagecolor}
\pagecolor{pagebg}
\color{bodyfg}

% ---- Fonts ------------------------------------------------------------
\usepackage[T1]{fontenc}
\IfFileExists{libertine.sty}{%
    \usepackage{libertine}      % warm, readable serif
    \usepackage[scaled=0.95]{inconsolata}  % softer monospace
}{}

% ---- Hyperref ---------------------------------------------------------
\usepackage{hyperref}
\hypersetup{
    colorlinks=true,
    linkcolor=linkcol,
    urlcolor=linkcol,
    citecolor=linkcol,
    pdfborder={0 0 0},
}

% ---- Section styling --------------------------------------------------
\usepackage[explicit]{titlesec}
\titleformat{\section}
    {\color{heading}\Large\bfseries}
    {\color{heading}\thesection\hspace{0.6em}}{0pt}{#1}
\titleformat{\subsection}
    {\color{subheading}\large\bfseries}
    {\color{subheading}\thesubsection\hspace{0.5em}}{0pt}{#1}
\titleformat{\subsubsection}
    {\color{accent}\normalsize\bfseries}
    {\color{accent}\thesubsubsection\hspace{0.5em}}{0pt}{#1}
\titleformat{\paragraph}[runin]
    {\color{accent}\bfseries}{}{0pt}{#1.}

\titlespacing*{\section}{0pt}{1.6em}{0.6em}
\titlespacing*{\subsection}{0pt}{1.2em}{0.4em}

% ---- Enumerate / itemize spacing --------------------------------------
\usepackage{enumitem}
\setlist[itemize]{leftmargin=1.4em, itemsep=2pt, topsep=4pt}
\setlist[enumerate]{leftmargin=1.6em, itemsep=2pt, topsep=4pt}

% ---- Code listings ----------------------------------------------------
\usepackage{listings}
\lstdefinestyle{notescode}{
    backgroundcolor=\color{codeBack},
    basicstyle=\ttfamily\small\color{bodyfg},
    keywordstyle=\color{codeKeyword}\bfseries,
    stringstyle=\color{codeString},
    commentstyle=\color{codeComment}\itshape,
    frame=single,
    framesep=6pt,
    rulecolor=\color{codeBorder},
    showstringspaces=false,
    breaklines=true,
    xleftmargin=10pt,
    xrightmargin=10pt,
    aboveskip=8pt,
    belowskip=8pt,
}
\lstset{style=notescode}

% ---- Callout boxes (tcolorbox) ----------------------------------------
\usepackage[most]{tcolorbox}

\newtcolorbox{defnbox}[1][]{
    enhanced, breakable,
    colback=defnBack, colframe=defnFrame, coltitle=bodyfg,
    title={\bfseries Definition \textemdash\ #1},
    fonttitle=\bfseries,
    boxrule=0.5pt, arc=2pt, left=8pt, right=8pt, top=6pt, bottom=6pt,
    attach boxed title to top left={xshift=10pt, yshift=-6pt},
    boxed title style={colback=defnFrame, colframe=defnFrame, arc=2pt},
    before skip=8pt, after skip=8pt,
}
\newtcolorbox{ideabox}[1][]{
    enhanced, breakable,
    colback=ideaBack, colframe=ideaFrame, coltitle=bodyfg,
    title={\bfseries Key idea #1},
    boxrule=0.5pt, arc=2pt, left=8pt, right=8pt, top=6pt, bottom=6pt,
    attach boxed title to top left={xshift=10pt, yshift=-6pt},
    boxed title style={colback=ideaFrame, colframe=ideaFrame, arc=2pt},
    before skip=8pt, after skip=8pt,
}
\newtcolorbox{warnbox}[1][]{
    enhanced, breakable,
    colback=warnBack, colframe=warnFrame, coltitle=bodyfg,
    title={\bfseries Gotcha #1},
    boxrule=0.5pt, arc=2pt, left=8pt, right=8pt, top=6pt, bottom=6pt,
    attach boxed title to top left={xshift=10pt, yshift=-6pt},
    boxed title style={colback=warnFrame, colframe=warnFrame, arc=2pt},
    before skip=8pt, after skip=8pt,
}
\newtcolorbox{notebox}[1][]{
    enhanced, breakable,
    colback=noteBack, colframe=noteFrame, coltitle=bodyfg,
    title={\bfseries Aside #1},
    boxrule=0.5pt, arc=2pt, left=8pt, right=8pt, top=6pt, bottom=6pt,
    attach boxed title to top left={xshift=10pt, yshift=-6pt},
    boxed title style={colback=noteFrame, colframe=noteFrame, arc=2pt},
    before skip=8pt, after skip=8pt,
}
\newtcolorbox{examplebox}[1][]{
    enhanced, breakable,
    colback=exBack, colframe=exFrame, coltitle=bodyfg,
    title={\bfseries Example #1},
    boxrule=0.5pt, arc=2pt, left=8pt, right=8pt, top=6pt, bottom=6pt,
    attach boxed title to top left={xshift=10pt, yshift=-6pt},
    boxed title style={colback=exFrame, colframe=exFrame, arc=2pt},
    before skip=8pt, after skip=8pt,
}

% ---- TikZ graph styles (added 2026-04-25 for ch_10 augmentation) ------
% tcolorbox[most] already loads tikz; this block declares the libraries
% + shared `vertex` / `edge` styles used by ch_10's general directed-graph
% diagrams. ch_9's tree-style TikZ uses inline overrides and is
% unaffected. Simple circle nodes + arrow edges only -- no production
% styling.
\usetikzlibrary{arrows.meta, positioning, calc}
\tikzset{
    vertex/.style={circle, draw, minimum size=7mm, inner sep=0pt,
                   font=\small, thick},
    vsource/.style={vertex, fill=ideaBack},
    vdone/.style={vertex, fill=defnBack},
    edge/.style={->, >=Stealth, thick},
    uedge/.style={-, thick},
    treeedge/.style={->, >=Stealth, thick},
    backedge/.style={->, >=Stealth, thick, dashed, red!70!black},
    fwdedge/.style={->, >=Stealth, thick, blue!60!black},
    crossedge/.style={->, >=Stealth, thick, green!50!black, dotted},
    mstedge/.style={-, very thick, accent},
}

% ---- Title block ------------------------------------------------------
\usepackage{titling}
\pretitle{\begin{center}\color{heading}\LARGE\bfseries}
\posttitle{\end{center}\vskip -0.4em}
\preauthor{\begin{center}\color{mutedtext}\normalsize}
\postauthor{\end{center}}
\predate{\begin{center}\color{mutedtext}\small}
\postdate{\end{center}\vspace{0.4em}{\color{rulecol}\hrule height 0.4pt}\vspace{1.2em}}
` right after \documentclass
%   in each chapter's lectures.tex and notes.tex.
% =======================================================================

% ---- Layout -----------------------------------------------------------
\usepackage[margin=1in]{geometry}
\usepackage{parskip}        % space between paragraphs, no indent
\usepackage{microtype}      % subtle kerning/spacing improvements

% ---- Colors (muted palette, easy on light-sensitive eyes) -------------
\usepackage{xcolor}
\definecolor{pagebg}       {HTML}{FAF6F0}  % warm cream background
\definecolor{bodyfg}       {HTML}{3D3A36}  % warm dark gray body text
\definecolor{heading}      {HTML}{3D6A6B}  % muted teal
\definecolor{subheading}   {HTML}{5B7A9C}  % dusty blue
\definecolor{accent}       {HTML}{8A6B4A}  % warm tan
\definecolor{linkcol}      {HTML}{5B7A9C}  % same as subheading
\definecolor{mutedtext}    {HTML}{6B6560}  % secondary / caption text
\definecolor{rulecol}      {HTML}{C8BFB0}  % separator / rule

% Callout box palette
\definecolor{defnFrame}    {HTML}{9DB89B}  \definecolor{defnBack}    {HTML}{EEF2E8}
\definecolor{ideaFrame}    {HTML}{9FB4CD}  \definecolor{ideaBack}    {HTML}{E8EEF5}
\definecolor{warnFrame}    {HTML}{CFAE87}  \definecolor{warnBack}    {HTML}{F5EDDF}
\definecolor{noteFrame}    {HTML}{BFA6AE}  \definecolor{noteBack}    {HTML}{F2E8EE}
\definecolor{exFrame}      {HTML}{9FBDBD}  \definecolor{exBack}      {HTML}{E8F0F0}
\definecolor{codeBack}     {HTML}{F3EDE2}
\definecolor{codeBorder}   {HTML}{D4CBBA}
\definecolor{codeKeyword}  {HTML}{6B4F7A}
\definecolor{codeString}   {HTML}{8A6B4A}
\definecolor{codeComment}  {HTML}{8A857F}

% ---- Page background + base text color --------------------------------
\usepackage{pagecolor}
\pagecolor{pagebg}
\color{bodyfg}

% ---- Fonts ------------------------------------------------------------
\usepackage[T1]{fontenc}
\IfFileExists{libertine.sty}{%
    \usepackage{libertine}      % warm, readable serif
    \usepackage[scaled=0.95]{inconsolata}  % softer monospace
}{}

% ---- Hyperref ---------------------------------------------------------
\usepackage{hyperref}
\hypersetup{
    colorlinks=true,
    linkcolor=linkcol,
    urlcolor=linkcol,
    citecolor=linkcol,
    pdfborder={0 0 0},
}

% ---- Section styling --------------------------------------------------
\usepackage[explicit]{titlesec}
\titleformat{\section}
    {\color{heading}\Large\bfseries}
    {\color{heading}\thesection\hspace{0.6em}}{0pt}{#1}
\titleformat{\subsection}
    {\color{subheading}\large\bfseries}
    {\color{subheading}\thesubsection\hspace{0.5em}}{0pt}{#1}
\titleformat{\subsubsection}
    {\color{accent}\normalsize\bfseries}
    {\color{accent}\thesubsubsection\hspace{0.5em}}{0pt}{#1}
\titleformat{\paragraph}[runin]
    {\color{accent}\bfseries}{}{0pt}{#1.}

\titlespacing*{\section}{0pt}{1.6em}{0.6em}
\titlespacing*{\subsection}{0pt}{1.2em}{0.4em}

% ---- Enumerate / itemize spacing --------------------------------------
\usepackage{enumitem}
\setlist[itemize]{leftmargin=1.4em, itemsep=2pt, topsep=4pt}
\setlist[enumerate]{leftmargin=1.6em, itemsep=2pt, topsep=4pt}

% ---- Code listings ----------------------------------------------------
\usepackage{listings}
\lstdefinestyle{notescode}{
    backgroundcolor=\color{codeBack},
    basicstyle=\ttfamily\small\color{bodyfg},
    keywordstyle=\color{codeKeyword}\bfseries,
    stringstyle=\color{codeString},
    commentstyle=\color{codeComment}\itshape,
    frame=single,
    framesep=6pt,
    rulecolor=\color{codeBorder},
    showstringspaces=false,
    breaklines=true,
    xleftmargin=10pt,
    xrightmargin=10pt,
    aboveskip=8pt,
    belowskip=8pt,
}
\lstset{style=notescode}

% ---- Callout boxes (tcolorbox) ----------------------------------------
\usepackage[most]{tcolorbox}

\newtcolorbox{defnbox}[1][]{
    enhanced, breakable,
    colback=defnBack, colframe=defnFrame, coltitle=bodyfg,
    title={\bfseries Definition \textemdash\ #1},
    fonttitle=\bfseries,
    boxrule=0.5pt, arc=2pt, left=8pt, right=8pt, top=6pt, bottom=6pt,
    attach boxed title to top left={xshift=10pt, yshift=-6pt},
    boxed title style={colback=defnFrame, colframe=defnFrame, arc=2pt},
    before skip=8pt, after skip=8pt,
}
\newtcolorbox{ideabox}[1][]{
    enhanced, breakable,
    colback=ideaBack, colframe=ideaFrame, coltitle=bodyfg,
    title={\bfseries Key idea #1},
    boxrule=0.5pt, arc=2pt, left=8pt, right=8pt, top=6pt, bottom=6pt,
    attach boxed title to top left={xshift=10pt, yshift=-6pt},
    boxed title style={colback=ideaFrame, colframe=ideaFrame, arc=2pt},
    before skip=8pt, after skip=8pt,
}
\newtcolorbox{warnbox}[1][]{
    enhanced, breakable,
    colback=warnBack, colframe=warnFrame, coltitle=bodyfg,
    title={\bfseries Gotcha #1},
    boxrule=0.5pt, arc=2pt, left=8pt, right=8pt, top=6pt, bottom=6pt,
    attach boxed title to top left={xshift=10pt, yshift=-6pt},
    boxed title style={colback=warnFrame, colframe=warnFrame, arc=2pt},
    before skip=8pt, after skip=8pt,
}
\newtcolorbox{notebox}[1][]{
    enhanced, breakable,
    colback=noteBack, colframe=noteFrame, coltitle=bodyfg,
    title={\bfseries Aside #1},
    boxrule=0.5pt, arc=2pt, left=8pt, right=8pt, top=6pt, bottom=6pt,
    attach boxed title to top left={xshift=10pt, yshift=-6pt},
    boxed title style={colback=noteFrame, colframe=noteFrame, arc=2pt},
    before skip=8pt, after skip=8pt,
}
\newtcolorbox{examplebox}[1][]{
    enhanced, breakable,
    colback=exBack, colframe=exFrame, coltitle=bodyfg,
    title={\bfseries Example #1},
    boxrule=0.5pt, arc=2pt, left=8pt, right=8pt, top=6pt, bottom=6pt,
    attach boxed title to top left={xshift=10pt, yshift=-6pt},
    boxed title style={colback=exFrame, colframe=exFrame, arc=2pt},
    before skip=8pt, after skip=8pt,
}

% ---- TikZ graph styles (added 2026-04-25 for ch_10 augmentation) ------
% tcolorbox[most] already loads tikz; this block declares the libraries
% + shared `vertex` / `edge` styles used by ch_10's general directed-graph
% diagrams. ch_9's tree-style TikZ uses inline overrides and is
% unaffected. Simple circle nodes + arrow edges only -- no production
% styling.
\usetikzlibrary{arrows.meta, positioning, calc}
\tikzset{
    vertex/.style={circle, draw, minimum size=7mm, inner sep=0pt,
                   font=\small, thick},
    vsource/.style={vertex, fill=ideaBack},
    vdone/.style={vertex, fill=defnBack},
    edge/.style={->, >=Stealth, thick},
    uedge/.style={-, thick},
    treeedge/.style={->, >=Stealth, thick},
    backedge/.style={->, >=Stealth, thick, dashed, red!70!black},
    fwdedge/.style={->, >=Stealth, thick, blue!60!black},
    crossedge/.style={->, >=Stealth, thick, green!50!black, dotted},
    mstedge/.style={-, very thick, accent},
}

% ---- Title block ------------------------------------------------------
\usepackage{titling}
\pretitle{\begin{center}\color{heading}\LARGE\bfseries}
\posttitle{\end{center}\vskip -0.4em}
\preauthor{\begin{center}\color{mutedtext}\normalsize}
\postauthor{\end{center}}
\predate{\begin{center}\color{mutedtext}\small}
\postdate{\end{center}\vspace{0.4em}{\color{rulecol}\hrule height 0.4pt}\vspace{1.2em}}


\title{CS 300 -- Chapter 2 Lectures\\\large Introduction to Algorithms}
\author{Jose de Lima}
\date{\today}

\begin{document}
\maketitle

\begin{ideabox}[Chapter map]
\textbf{Where this sits in CS 300.} Chapter 1 taught you how to \emph{type}
a loop; chapter 2 teaches you how to \emph{decide which loop to write} when
you haven't been given one. This is the bridge between ``I can code'' and
``I can solve a novel problem.'' Every later chapter is a specific data
structure; this chapter is the \emph{meta-skill} of picking an algorithmic
strategy in the first place.

\textbf{What you add to your toolkit here:}
\begin{itemize}
    \item A vocabulary for problems vs.\ algorithms, and for what makes
          one algorithm ``better'' than another (correctness, speed,
          memory, simplicity).
    \item A \textbf{recursive mindset}: base case + smaller subproblem.
          Factorial and binary search are the training wheels; BSTs,
          merge sort, and graph traversal all run on this.
    \item Three named strategies you'll reach for repeatedly:
          \textbf{heuristic}, \textbf{greedy}, \textbf{dynamic programming}.
          Knowing when each is appropriate is half the exam.
    \item An ethics / data-privacy framework so that the code you write
          doesn't just \emph{work} but is defensible when someone asks
          why you chose to collect or expose certain data.
\end{itemize}

\textbf{Mastery checklist} -- you should be able to, without references:
\begin{enumerate}
    \item State the two ingredients every recursive function needs,
          and show what goes wrong if either is missing (infinite
          recursion / wrong answer).
    \item Write \texttt{factorial(n)} and \texttt{binarySearch(v, t)}
          recursively in C++ \emph{and} iteratively, and argue when each
          form is appropriate.
    \item Explain why na\"ive recursive Fibonacci is exponential while
          the iterative / memoized version is linear -- in terms of the
          recursion tree, not hand-waving.
    \item Given a new problem, correctly classify it as a candidate for
          greedy, DP, or neither, and defend the choice with one of the
          structural tests (greedy-choice property, optimal substructure,
          overlapping subproblems).
    \item Walk through the fractional-knapsack greedy proof (``always
          pick highest value/weight ratio next'') and know why the
          \emph{0/1} variant is not solvable greedily.
    \item Identify overlapping subproblems in an LCS-like problem and
          write the 2D recurrence relation before touching code.
    \item Name at least one privacy or ethics consideration that could
          change your design decision on a data-driven feature.
\end{enumerate}

\textbf{Looking ahead.} Chapter 3 gives you the formal language (Big-O)
to justify the ``faster'' and ``slower'' claims you made informally here.
The sorting algorithms in ch.~3 are concrete instances of the strategies
here: insertion/selection sort are incremental, merge/quicksort are
divide-and-conquer, radix sort is a non-comparison heuristic. Expect ch.~3
to be easier if your ch.~2 recursion intuitions are solid.
\end{ideabox}

\section{2.1 Introduction to Algorithms}

Chapter 1 covered the \emph{mechanics} of C++ -- vectors, strings,
loops. Chapter 2 starts asking a different question: given a problem,
what sequence of steps actually solves it, and how do you compare one
sequence against another? That is the subject of \textbf{algorithms},
and it is the real content of this course.

\subsection{Problems vs. algorithms}

Two words that get used interchangeably but mean different things.

\begin{defnbox}[Computational problem]
A \textbf{computational problem} is a precise specification of three
things: the \textbf{input} the solver will receive, the
\textbf{question} being asked about that input, and the \textbf{output}
the solver should produce. It describes \emph{what} needs to be done,
not \emph{how}.
\end{defnbox}

\begin{defnbox}[Algorithm]
An \textbf{algorithm} is a finite, unambiguous sequence of steps that
solves a computational problem. Algorithms describe \emph{how} the
problem gets solved. The same problem can have many algorithms; the
interesting question is usually which one is best.
\end{defnbox}

\begin{examplebox}[A problem and one algorithm that solves it]
\textbf{Problem:} Given an array of numbers, what is the maximum value?

\textbf{Algorithm (pseudocode):}
\begin{verbatim}
FindMax(a):
    max = a[0]
    for i = 1 to a.size - 1:
        if a[i] > max:
            max = a[i]
    return max
\end{verbatim}
The problem specification does not mention a loop, or
\texttt{max}, or \texttt{i}. Those are \emph{implementation decisions}
the algorithm makes. A different algorithm -- e.g., sort the array and
return the last element -- solves the same problem with different
steps and different efficiency.
\end{examplebox}

\begin{notebox}[Algorithms exist independent of language]
Algorithms are usually described in pseudocode, diagrams, or English
because the language is incidental to the \emph{idea}. The C++
implementation you will write in homework is just one translation of
the algorithm. Binary search is binary search whether it is written
in C++, Python, or executed by a human with a pencil.
\end{notebox}

\subsection{Why this matters: picking the right algorithm is the job}

Most interesting problems you will ever see have already been reduced
to a small set of well-known subproblems, each with a well-studied
algorithm. Half of being a good engineer is recognizing which known
problem the thing in front of you actually is.

\begin{center}
\small
\begin{tabular}{lll}
\textbf{Domain} & \textbf{Problem} & \textbf{Common algorithm} \\
\hline
DNA analysis & Longest shared subsequence & Longest common substring \\
E-commerce / search & ``Is this product in stock?'' & Binary search \\
Navigation & Fastest route A $\to$ B & Dijkstra's shortest path \\
Scheduling & Which tasks can run first? & Topological sort \\
Compilers / build & What to rebuild after a change & DFS on a dependency graph \\
Spell check / autocomplete & Closest word to a typo & Edit distance (DP) \\
\end{tabular}
\normalsize
\end{center}

The CS 300 syllabus is essentially a tour of the right-hand column --
sorting, searching, hashing, trees, heaps, graph traversal. Once you
can name the tool, the implementation becomes mechanical.

\subsection{What makes an algorithm ``good''?}

The dominant metric is \textbf{runtime growth} -- how the number of
steps changes as the input gets bigger. This is the subject of Chapter
2 proper.

\begin{defnbox}[Efficient algorithm (informal)]
An algorithm is \textbf{efficient} if its runtime grows at most
\textbf{polynomially} with the input size: $n$, $n \log n$, $n^2$,
$n^3$, etc. Anything exponential ($2^n$, $n!$) is considered too slow
to scale.
\end{defnbox}

\begin{notebox}[Why polynomial is the cutoff]
It is an arbitrary line, but a useful one. Polynomial-time algorithms
remain tractable as inputs grow -- doubling $n$ multiplies the runtime
by a constant factor. Exponential algorithms do not -- doubling $n$
can \emph{square} the runtime. On a modern machine, $2^{40}$ is
measured in minutes; $2^{60}$ is measured in years. That is the
difference between ``run it overnight'' and ``never finishes.''
\end{notebox}

\subsection{NP-complete problems: when no fast algorithm is known}

Some problems have resisted every attempt to find an efficient
algorithm. These are the \textbf{NP-complete} problems.

\begin{defnbox}[NP-complete (intuitive)]
A problem is \textbf{NP-complete} if:
\begin{itemize}[leftmargin=*,noitemsep]
    \item nobody has found a polynomial-time algorithm for it,
    \item nobody has proven that one is impossible, and
    \item a fast algorithm for \emph{any} NP-complete problem would
          give a fast algorithm for \emph{all} of them.
\end{itemize}
Whether \textbf{P = NP} -- i.e., whether such an algorithm exists -- is
the most famous open problem in computer science.
\end{defnbox}

\begin{examplebox}[The clique problem]
\textbf{Input:} a social-network graph and an integer $k$. \\
\textbf{Question:} is there a set of $k$ people who all know each other?

Finding a clique of size 3 in a small graph is easy. Finding one of
size 50 in a graph of millions of nodes is, as far as anyone knows,
beyond any algorithm that could finish this century. Same problem
statement, exponentially harder instance.
\end{examplebox}

\begin{ideabox}[What you do when a problem is NP-complete]
Not ``give up.'' The usual response is to solve a \emph{relaxed}
version:
\begin{itemize}[leftmargin=*,noitemsep]
    \item \textbf{Approximation:} find a provably-close-to-optimal
          answer quickly.
    \item \textbf{Heuristic:} use a good-enough rule of thumb that
          works on typical inputs.
    \item \textbf{Special-case:} if the input has extra structure
          (sparse graph, small $k$, ordered input), a faster algorithm
          may exist for that subset.
\end{itemize}
Recognizing a problem as NP-complete saves you from hunting for an
efficient exact algorithm that doesn't exist. That is a skill in
itself.
\end{ideabox}

\begin{ideabox}[Big picture]
This course teaches two things simultaneously: a catalog of named
algorithms that handle the most common computational problems, and a
\emph{language} for comparing algorithms (big-O, which is formalized
in chapter~3) so you can justify your choice. When you finish, the
goal is not that you have memorized every pseudocode listing, but
that you can look at a new problem, say ``that's a search on sorted
data, so binary search is O($\log n$)'' -- and pick up the right tool
without rediscovering it.
\end{ideabox}

\section{2.2 Problem Solving}

Knowing C++ is a prerequisite for writing code, but it is not what
makes programs \emph{work}. The hard part -- the part people are
actually paid for -- is \textbf{problem solving}: figuring out a
methodical sequence of steps that turns the input into the desired
output. The programming language is just the notation you write that
sequence down in, the way English is the notation a chef writes a
recipe down in. Good recipes come from good cooking, not good English.

\begin{defnbox}[Solution approach]
A \textbf{solution approach} is the idea of how you will solve a
problem, stated in words, diagrams, or pseudocode -- independent of
any programming language. A good solution approach is worked out
before any code is written; the code is the final translation, not
the design.
\end{defnbox}

\subsection{A non-programming example: matching socks}

Three kinds of socks sit unsorted in a drawer. You want one matching
pair.

\begin{examplebox}[Two approaches, very different runtimes]
\textbf{Approach A -- ``try and put back'':} pull a sock, pull another,
compare. If they match, done. Otherwise put the second one back and
try again.

\textbf{Approach B -- ``pull and keep'':} pull a sock, then another,
then another, stopping only when you have a pair.

Approach A has no guaranteed bound -- you could get unlucky for a
long time. Approach B is bounded: with $k$ sock types, pulling
$k + 1$ socks \emph{must} produce a match (this is the pigeonhole
principle). For 3 sock types, four pulls is enough, always.
\end{examplebox}

\begin{notebox}[The pigeonhole principle]
If $n+1$ items go into $n$ boxes, some box must contain at least
two items. This comes back repeatedly -- hash-table collisions, proofs
about sorting lower bounds, birthday-paradox arguments -- so
recognizing it early is worth it. In the sock problem, the ``boxes''
are the sock types and the ``items'' are the socks you've pulled.
\end{notebox}

\subsection{A programming-flavored example: sorting name tags}

1000 name tags, currently sorted by first name, need to be re-sorted
by last name.

\begin{examplebox}[Two approaches]
\textbf{Approach 1 -- ``insert each into a new sorted stack'':}
take one tag at a time, walk down a new stack, drop it in the right
place. Roughly $n^2$ comparisons, because each insertion scans
everything already placed.

\textbf{Approach 2 -- ``bucket then sort each bucket'':}
make 26 sub-stacks, one per starting letter of the last name. Drop
each tag into its letter's bucket. Sort each bucket individually.
Concatenate the buckets in order.
\end{examplebox}

\begin{ideabox}[This is actually radix sort in disguise]
Approach 2 is a hand-executed version of \textbf{radix sort} --
partition by one ``digit'' of the key (here, the first letter of the
last name), sort within each partition, and recombine. You will meet
radix sort formally later in the course. The point here is that the
two approaches give the same answer but the work is distributed very
differently, and one of them is a named algorithm with well-known
properties. Almost every ``clever'' approach you invent on a
whiteboard turns out to be a known algorithm.
\end{ideabox}

\subsection{The 64-person greeting problem}

An organizer wants every person at a 64-person meeting to greet every
other person for 30 seconds. How long does this take?

\begin{examplebox}[Doing the arithmetic]
64 people, each pair meets once, so the number of pairs is
$\binom{64}{2} = \frac{64 \cdot 63}{2} = 2016$ pairs. At 30 seconds
each, that is $2016 \cdot 30 = 60{,}480$ seconds $\approx 16.8$ hours.
\end{examplebox}

\begin{ideabox}[The combinatorial trap]
Schemes of the form ``everybody does something with everybody else''
are $O(n^2)$. They \emph{look} linear in the problem statement
(``just have everyone say hi'') but scale quadratically in practice.
This is the same reason all-pairs-shortest-path takes $O(n^3)$, why
bubble sort is $O(n^2)$, and why naively checking every pair of
records in a database does not work at a million-row scale. Watch
for ``every X with every other X'' in a problem statement -- it's
a flag for $O(n^2)$.
\end{ideabox}

\subsection{A usable workflow}

The approach this course will try to instill, roughly in order:

\begin{enumerate}[leftmargin=*]
    \item \textbf{Read the problem until you can restate it in one
          sentence} without looking at the prompt. If you cannot, you
          haven't understood it yet; coding will not help.
    \item \textbf{Solve a tiny example by hand.} Before writing code,
          walk through the problem with a three- or four-element
          input. This usually exposes missing details and suggests
          an algorithm.
    \item \textbf{Describe an approach in pseudocode or English.}
          Steps, not syntax. If the approach doesn't work on paper,
          C++ will not save it.
    \item \textbf{Estimate the cost.} Will this scale to the largest
          input the problem allows? If no, find a better approach
          before you write any code.
    \item \textbf{Translate to C++.} Now the language matters.
          Idiomatic vectors, bounds checking, \texttt{.at()} while
          debugging, algorithms from \texttt{<algorithm>} where they
          fit.
    \item \textbf{Test with the tiny example from step 2.} If the
          output does not match, trace the loop (\textsection 1.10
          methodology).
\end{enumerate}

\begin{warnbox}[The anti-pattern: code first, think later]
The most common beginner mistake is to start typing at step 5, work
out the approach through compiler errors, and debug by randomly
changing conditions. It \emph{feels} productive (there are always
more changes to try) but often results in code that passes the
examples you happened to test and fails everything else. The
up-front thinking from steps 1--4 takes minutes and saves hours.
\end{warnbox}

\begin{notebox}[Two paths to an algorithm: reduce, or design]
The workflow above leaves the question of \emph{which} algorithm
implicit. MIT 6.006 frames it as a top-level binary choice:

\begin{enumerate}[leftmargin=*,noitemsep]
    \item \textbf{Reduce to a problem you already know.} If the
          problem looks like \emph{search}, reach for a search data
          structure (array, hash table, BST). If \emph{sort}, reach
          for a sort algorithm (insertion, merge, radix). If
          \emph{shortest path}, a graph algorithm (BFS, Dijkstra,
          Bellman-Ford). The cs-300 syllabus is essentially this
          catalog of named problems and their canonical solutions.
    \item \textbf{Design your own (recursive) algorithm.} If no
          known problem fits, you're inventing. The strategies in
          \textsection 2.4--\textsection 2.10 (recursion,
          divide-and-conquer, DP, greedy, heuristics) are the
          design patterns for that case.
\end{enumerate}

The reduce path is faster when it applies, which is most of the
time on a coursework or job-interview problem. Train the
recognition reflex first; the design strategies are what you fall
back on when recognition fails.
\end{notebox}

\begin{ideabox}[Big picture]
A programming language is a notation for expressing solutions; it is
not itself a solution. The thing you develop in this course is the
habit of looking at a problem, recognizing a structure you have seen
before, picking the matching algorithm, and only then writing the
C++. When you get stuck on a homework problem, the fix is almost
never ``more C++'' -- it is going back to pseudocode or walking
through a 3-element input by hand.
\end{ideabox}

\section{2.3 Computational Problem Solving}

\textsection 2.2 framed problem solving in general terms. This section
narrows in on the specific style of thinking that computing demands,
which cognitive-science and CS-education literature has settled on
calling \textbf{computational thinking}. It is four habits --
decomposition, pattern recognition, algorithms, and abstraction --
applied together. The names are generic; the habits are not.

\subsection{The input > process > output model}

Every program, no matter how large, fits the same shape at the top
level: take input, compute, produce output (and/or persistent side
effects). Keeping this shape in mind while designing prevents the
common trap of conflating ``how data gets in'' with ``what we do with
it.''

\begin{examplebox}[Primality check, top-level view]
\begin{verbatim}
INPUT:    a positive integer n
PROCESS:  decide whether n has any divisors other than 1 and n
OUTPUT:   "prime" or "not prime"
\end{verbatim}
The process step is where the algorithm lives. The input and output
steps are routine -- read an int, print a string -- and should not
dominate your design time.
\end{examplebox}

\subsection{The four habits of computational thinking}

\begin{defnbox}[Computational thinking]
A problem-solving approach built from four interacting moves:
\begin{itemize}[leftmargin=*,noitemsep]
    \item \textbf{Decomposition} -- break a big problem into smaller
          pieces.
    \item \textbf{Pattern recognition} -- notice repeated structure
          within or across problems.
    \item \textbf{Algorithms} -- express the solution as a precise
          sequence of steps.
    \item \textbf{Abstraction} -- generalize so one solution handles a
          whole family of problems.
\end{itemize}
These are not stages in a waterfall. They feed each other: decomposing
reveals patterns, patterns suggest algorithms, a finished algorithm
gets abstracted so it can be reused.
\end{defnbox}

\subsection{Worked example: is $n$ prime?}

Watch the four habits in action on one problem.

\paragraph{Decomposition.} Big question (``is $n$ prime?'') into
smaller yes/no questions: is it divisible by 2? by 3? by 4? \ldots\ If
any answer is yes, $n$ is not prime; if all are no, it is prime.

\paragraph{Pattern recognition.} Two useful observations:
\begin{itemize}[leftmargin=*,noitemsep]
    \item We can stop the first time we see a divisor -- we don't need
          to check the rest.
    \item No factor of $n$ (other than $n$ itself) can exceed $n/2$,
          so we never need to test past there. In fact, no factor
          exceeds $\sqrt{n}$, which is a much tighter bound.
\end{itemize}

\paragraph{Algorithm.}
\begin{verbatim}
IsPrime(n):
    if n < 2: return false
    for d = 2 to sqrt(n):
        if n mod d == 0: return false
    return true
\end{verbatim}

\paragraph{Abstraction.} The algorithm doesn't depend on any specific
value of $n$; it works for every positive integer. Expressed as a
function it becomes a reusable tool.

\begin{examplebox}[C++ translation]
\begin{verbatim}
#include <cmath>

bool is_prime(int n) {
    if (n < 2) return false;
    int limit = static_cast<int>(std::sqrt(n));
    for (int d = 2; d <= limit; ++d) {
        if (n % d == 0) return false;
    }
    return true;
}
\end{verbatim}
Three pattern-recognition wins compound here: stop at \texttt{sqrt(n)}
instead of $n/2$, return as soon as a divisor is found, and handle the
$n < 2$ edge case up front so the loop body doesn't have to.
\end{examplebox}

\begin{notebox}[Why $\sqrt{n}$ is enough]
If $n = a \cdot b$ with $a \le b$, then $a \le \sqrt{n}$. So if $n$
has any factor at all, it has one no bigger than $\sqrt{n}$. Checking
past $\sqrt{n}$ is redundant. Turning a $O(n)$ primality check into
$O(\sqrt{n})$ is not a rewrite; it's a single-line loop bound change
that pattern-recognition hands you for free.
\end{notebox}

\subsection{Second example: sum of 1 to 100}

\paragraph{Decomposition.} Add numbers one at a time:
$1 + 2$, then $+ 3$, \ldots

\paragraph{Pattern recognition.} Pair the numbers from the ends in:
$(1 + 100), (2 + 99), (3 + 98), \ldots$ Each pair sums to 101, and
there are 50 pairs. Total: $50 \cdot 101 = 5050$.

\paragraph{Algorithm.} The closed-form $\dfrac{n(n+1)}{2}$ computes
the sum in $O(1)$ rather than $O(n)$.

\paragraph{Abstraction.} The formula works for any positive $n$, not
just 100. It generalizes further to arithmetic sequences.

\begin{examplebox}[Two implementations, same output]
\begin{verbatim}
// O(n) loop -- decomposition only
int sum_loop(int n) {
    int s = 0;
    for (int i = 1; i <= n; ++i) s += i;
    return s;
}

// O(1) closed form -- pattern recognition plus algebra
int sum_formula(int n) {
    return n * (n + 1) / 2;
}
\end{verbatim}
For $n = 100$ both return 5050. For $n = 1{,}000{,}000{,}000$, only
the second one finishes in microseconds.
\end{examplebox}

\begin{ideabox}[Pattern recognition is the most valuable habit]
Of the four, pattern recognition is where the real leverage lives.
Decomposition and algorithms are mechanical once the pattern is
clear; abstraction is mostly cleanup. Spotting that ``this is a
sum of pairs'' or ``this is a divisibility scan'' turns a tedious
loop into a one-line expression, and recognizing ``this is binary
search'' turns an $O(n)$ scan into $O(\log n)$. Invest in the
vocabulary of named patterns -- this course is partly that vocabulary.
\end{ideabox}

\subsection{The loop from \textsection 2.2 revisited}

The workflow from the previous section is computational thinking in
practice:

\begin{center}
\small
\begin{tabular}{ll}
\textbf{Workflow step} & \textbf{CT concept} \\
\hline
Restate the problem & (prerequisite) \\
Solve a tiny example by hand & decomposition \\
Spot the structure & pattern recognition \\
Write pseudocode & algorithm \\
Generalize to a function & abstraction \\
Estimate the cost & pattern recognition (again) \\
Translate to C++ & implementation \\
Test on the tiny example & verification \\
\end{tabular}
\normalsize
\end{center}

\begin{ideabox}[Big picture]
Computational thinking is not a separate subject from programming; it
is the shape of the thinking programming rewards. Every named
algorithm in this course -- binary search, merge sort, Dijkstra, hash
tables -- is someone's decomposition-plus-pattern-recognition written
down and abstracted. The more of them you know, the less original
design you need to do on any given problem: you recognize the shape,
retrieve the matching algorithm, and translate it. That is what makes
experienced engineers productive.
\end{ideabox}

\section{2.4 Recursive Definitions}

A \textbf{recursive} algorithm solves a problem by reducing it to a
\emph{smaller instance of the same problem} and combining the answer.
It is the same mental move as mathematical induction, and it unlocks a
surprising number of problems that are awkward to express with loops
alone -- tree traversals, divide-and-conquer sorts, graph searches,
backtracking puzzles. The entire back half of this course leans on it.

\begin{defnbox}[Recursive algorithm]
An algorithm is \textbf{recursive} if it solves a problem by (1)
identifying one or more \textbf{base cases} that can be answered
directly, and (2) reducing every other input to a \emph{smaller}
instance of the same problem, calling itself on the smaller instance,
and combining the result.
\end{defnbox}

\begin{defnbox}[Base case]
The \textbf{base case} is the one (or few) input(s) that the algorithm
answers without recursing. Every recursive algorithm needs at least
one. Without it, the algorithm keeps reducing forever -- in code, the
call stack grows until the program crashes.
\end{defnbox}

\begin{defnbox}[Recursive case]
The \textbf{recursive case} is the work the algorithm does on a
non-base input: reduce the problem, call itself on the smaller
version, combine the returned answer with whatever local work
remains.
\end{defnbox}

\subsection{The anatomy of every recursive function}

\begin{examplebox}[The template]
\begin{verbatim}
ReturnType solve(Input x) {
    if (is_base_case(x)) {
        return direct_answer(x);        // base case
    }
    Input smaller = reduce(x);          // make progress
    ReturnType sub = solve(smaller);    // recurse
    return combine(x, sub);             // combine and return
}
\end{verbatim}
Every recursive function you will write is some fill-in-the-blank
version of this skeleton.
\end{examplebox}

\begin{defnbox}[SRTBOT: a fuller version of the template]
The 3-part template above (\emph{base, reduce, combine}) is the
core of recursive algorithm design. MIT 6.006 expands it into a
6-step recipe -- mnemonic \textbf{SRTBOT} -- that adds the explicit
checks the abbreviated template leaves implicit:

\begin{enumerate}[leftmargin=*,noitemsep]
    \item \textbf{S}ubproblem definition. State the meaning of a
          subproblem in words, in terms of parameters. Often
          subsets of input: prefixes, suffixes, contiguous
          substrings.
    \item \textbf{R}elate. Express the subproblem solution in terms
          of smaller subproblems. The design pattern: \emph{identify
          a question whose answer would reduce to smaller subproblems,
          then locally brute-force all possible answers}.
    \item \textbf{T}opological order. Argue the relation is
          acyclic -- subproblems form a DAG -- so the recursion
          terminates.
    \item \textbf{B}ase cases. Solutions for the smallest
          subproblems where the relation breaks down.
    \item \textbf{O}riginal problem. Show how to compute the answer
          to the original problem from subproblem solutions. Use
          parent pointers to recover the actual solution, not just
          the objective value.
    \item \textbf{T}ime analysis. Sum work over subproblems
          ($\sum_{x \in X} \text{work}(x)$, treating recursive
          calls as $O(1)$).
\end{enumerate}

The 3-part template is the SRTBOT core: \emph{base} = step 4,
\emph{reduce + combine} = step 2. SRTBOT pays off when you start
writing dynamic-programming algorithms in \textsection 2.10, where
the topological-order argument and the work-per-subproblem analysis
become the actual content of the design.
\end{defnbox}

\subsection{Example 1: factorial}

$n! = n \cdot (n-1) \cdot (n-2) \cdots 2 \cdot 1$, with $1! = 1$ as the
base.

\begin{examplebox}[Factorial in C++]
\begin{verbatim}
long long factorial(int n) {
    if (n <= 1) return 1;             // base case
    return n * factorial(n - 1);      // recursive case
}
\end{verbatim}
Trace for $n = 4$:
\begin{verbatim}
factorial(4) = 4 * factorial(3)
             = 4 * 3 * factorial(2)
             = 4 * 3 * 2 * factorial(1)
             = 4 * 3 * 2 * 1
             = 24
\end{verbatim}
\end{examplebox}

\begin{warnbox}[\texttt{n!} overflows 64-bit \texttt{long long} at $n = 21$]
$20! \approx 2.4 \times 10^{18}$, just under the $\texttt{long long}$
max. $21!$ silently wraps around to a nonsense value. For bigger
inputs you need a big-integer type or modular arithmetic. This isn't
a recursion bug -- iterative factorial has the same problem -- but
students often think the recursive version ``broke'' when the
arithmetic did.
\end{warnbox}

\subsection{Example 2: cumulative sum}

Sum of $0 + 1 + 2 + \cdots + n$. Base case: sum to $0$ is $0$.

\begin{examplebox}[\texttt{sum(n)}]
\begin{verbatim}
int cumulative_sum(int n) {
    if (n == 0) return 0;
    return n + cumulative_sum(n - 1);
}
\end{verbatim}
$\texttt{sum}(4) = 4 + 3 + 2 + 1 + 0 = 10$.

This is a great ``first recursion'' exercise, but the closed form
$n(n+1)/2$ is O(1) and doesn't risk stack overflow. Recursion here is
pedagogy, not the right tool.
\end{examplebox}

\subsection{Example 3: reverse a vector by swapping from both ends}

A recursive rewrite of the reversal loop from \textsection 1.9.

\begin{examplebox}[Reverse in place, recursively]
\begin{verbatim}
void reverse_list(std::vector<int>& v, int lo, int hi) {
    if (lo >= hi) return;              // base: 0 or 1 element left
    std::swap(v[lo], v[hi]);           // work at this layer
    reverse_list(v, lo + 1, hi - 1);   // recurse on the inside
}

// Caller:
reverse_list(v, 0, v.size() - 1);
\end{verbatim}
Compare to the loop in \textsection 1.9. The recursion mirrors the
structural invariant directly: \emph{the outer layer is done; reverse
the inner layer.} Some programmers find that clearer than the
\texttt{v.size() / 2} loop bound.
\end{examplebox}

\subsection{What can go wrong}

\begin{warnbox}[1. Missing or unreachable base case]
\begin{verbatim}
int factorial(int n) {
    return n * factorial(n - 1);   // never stops
}
\end{verbatim}
\texttt{factorial(0)} calls \texttt{factorial(-1)} calls
\texttt{factorial(-2)} \ldots\ the program crashes with a stack
overflow. Every recursive function needs a base case that's
\emph{reachable} from every legal input.
\end{warnbox}

\begin{warnbox}[2. Recursing on the same or larger subproblem]
\begin{verbatim}
int f(int n) {
    if (n == 0) return 0;
    return f(n) + 1;              // same n -- infinite recursion
}
\end{verbatim}
If the recursive call isn't on a \emph{strictly smaller} input (or
strictly closer to the base case), recursion doesn't terminate. A
useful habit: write down a \emph{measure} -- an integer that decreases
on every call -- and confirm it reaches the base.
\end{warnbox}

\begin{warnbox}[3. Stack overflow on deep recursion]
C++ call stacks are usually measured in megabytes; each stack frame
is a few dozen to a few hundred bytes. That makes the practical
recursion depth somewhere in the tens of thousands to millions of
frames before the stack blows. For chains like \texttt{factorial(n)}
or recursive sum this rarely matters, but for tree traversal on a
badly-shaped tree -- or for code that recurses on linked list length
-- you can hit it. Iteration or explicit stack data structures are
the fixes.
\end{warnbox}

\begin{notebox}[Recursive strategies, classified by call-graph shape]
A recursive call graph (vertex per call, directed edge $A \to B$ if
$B$ calls $A$) must be acyclic for the recursion to terminate. The
\emph{shape} of that DAG corresponds 1:1 with named algorithm-design
strategies you'll meet in \textsection 2.5--\textsection 2.10:

\begin{center}
\small
\begin{tabular}{ll}
\textbf{Strategy} & \textbf{Call-graph shape} \\
\hline
Brute force                & Star (all calls at one level) \\
Decrease \& conquer        & Chain (each call has one child) \\
Divide \& conquer          & Tree (calls split, no overlap) \\
Dynamic programming        & DAG (calls overlap; reuse via memo) \\
Greedy / incremental       & Subgraph (only one path explored) \\
\end{tabular}
\normalsize
\end{center}

This is the unifying framing for the rest of the chapter.
\textsection 2.5 binary search is decrease-and-conquer (chain);
naive Fibonacci is divide-and-conquer with overlap (a tree that
\emph{wants} to be a DAG); \textsection 2.10 DP is the corresponding
DAG fix; \textsection 2.9 greedy commits to one branch instead of
exploring the whole shape. When you can sketch the call graph for a
small input, you can usually name the strategy by inspection.
\end{notebox}

\subsection{When to reach for recursion}

\begin{ideabox}[Recursion shines on recursive structure]
Some data has recursion baked in -- \textbf{trees} have subtrees,
\textbf{lists} have tails, \textbf{grids} have sub-grids, \textbf{graph
searches} explore neighbors whose structure is ``the same problem on
a smaller set of nodes.'' When the data is recursive, recursive code
is shorter and clearer than the iterative equivalent. When the data
is linear and the work is constant-per-element, loops almost always
win.
\end{ideabox}

\subsection{Recursion and induction -- the same idea}

A correctness argument for recursion mirrors a proof by induction:

\begin{enumerate}[leftmargin=*,noitemsep]
    \item \textbf{Base case.} The base-case answer is correct by
          direct inspection.
    \item \textbf{Inductive step.} Assume the recursive call returns
          the correct answer for the smaller input. Show that the
          combine step produces the correct answer for the current
          input.
\end{enumerate}

\begin{examplebox}[Correctness of \texttt{factorial}]
Base: \texttt{factorial(1) == 1}, matching $1! = 1$. \\
Step: Assume \texttt{factorial(n-1) == (n-1)!}. Then
\texttt{n * factorial(n-1) == n * (n-1)! == n!}. By induction,
\texttt{factorial(n)} is correct for all $n \ge 1$.
\end{examplebox}

\begin{ideabox}[Big picture]
Recursion is worth the investment even when a loop would work,
because half of the algorithms in the rest of this course are most
naturally expressed recursively: binary search, merge sort, quicksort,
tree traversal, DFS, backtracking, dynamic programming recurrences.
Get comfortable with the three-part template (base, reduce, combine)
and the induction-style correctness argument, and those algorithms
stop feeling mysterious -- they are the same pattern applied to
different kinds of ``smaller subproblem.''
\end{ideabox}

\section{2.5 Recursive Algorithms}

\textsection 2.4 introduced the mechanics of recursion on small
arithmetic examples. This section applies the same three-part template
(base, reduce, combine) to two algorithms you will actually use:
\textbf{Fibonacci} (a warning sign about naive recursion) and
\textbf{binary search} (a textbook win for recursion). Read them side
by side -- same shape, wildly different efficiency.

\subsection{Fibonacci: the cautionary tale}

\begin{defnbox}[Fibonacci sequence]
$F_0 = 0$, $F_1 = 1$, and $F_n = F_{n-1} + F_{n-2}$ for $n \ge 2$. The
sequence begins $0, 1, 1, 2, 3, 5, 8, 13, 21, 34, \ldots$
\end{defnbox}

The definition \emph{is} a recurrence, so the recursive implementation
writes itself:

\begin{examplebox}[Naive recursive Fibonacci]
\begin{verbatim}
long long fib(int n) {
    if (n == 0) return 0;
    if (n == 1) return 1;
    return fib(n - 1) + fib(n - 2);
}
\end{verbatim}
Base cases: $n = 0$ and $n = 1$. Recursive case: sum of two smaller
subproblems.
\end{examplebox}

\begin{warnbox}[This function is catastrophically slow]
\texttt{fib(50)} takes minutes. \texttt{fib(60)} won't finish today.
The reason: \texttt{fib(n)} calls both \texttt{fib(n-1)} and
\texttt{fib(n-2)}, which each call two more, and so on. The number of
calls grows roughly like $\phi^n \approx 1.618^n$. The recursion tree
has billions of nodes, almost all of them computing values you've
already computed:
\begin{verbatim}
                fib(5)
               /      \
          fib(4)       fib(3)
         /     \      /     \
      fib(3)  fib(2) fib(2) fib(1)
      ...
\end{verbatim}
\texttt{fib(3)} is computed twice, \texttt{fib(2)} three times,
\texttt{fib(1)} five times. This is the canonical example of a
\emph{correct} recursive algorithm that is \emph{unacceptably slow}.
\end{warnbox}

\begin{ideabox}[Two standard fixes]
\textbf{Memoization:} remember each result the first time you compute
it. Recursion plus a cache turns $O(\phi^n)$ into $O(n)$.

\textbf{Iteration (bottom-up):} just build the sequence from the
bottom in a loop -- no recursion, $O(n)$ time, $O(1)$ space.
\begin{verbatim}
long long fib_iter(int n) {
    long long a = 0, b = 1;
    for (int i = 0; i < n; ++i) {
        long long next = a + b;
        a = b;
        b = next;
    }
    return a;
}
\end{verbatim}
The \emph{shape} of the recurrence makes naive recursion tempting.
Recognizing overlapping subproblems -- a job for pattern recognition
(\textsection 2.3) -- tells you to cache or iterate instead.
\end{ideabox}

\begin{notebox}[This is a preview of dynamic programming]
``Recursion with overlapping subproblems'' is the single-sentence
definition of \textbf{dynamic programming}. Every DP problem is a
recursive definition plus a cache (memoization) or a loop (tabulation).
You will meet this idea formally later; the Fibonacci blow-up is
where the need for it first becomes obvious.
\end{notebox}

\subsection{Binary search: the textbook win}

Here recursion earns its keep. Binary search on a \emph{sorted} array
finds a key in $O(\log n)$ comparisons -- dramatically faster than the
linear $O(n)$ scan from \textsection 1.5.

\begin{defnbox}[Binary search]
To search a sorted array for a key: look at the middle element. If it
equals the key, return. If the key is smaller, recurse on the left
half; if larger, recurse on the right half. If the range becomes
empty, the key is not present.
\end{defnbox}

\begin{examplebox}[Recursive binary search]
\begin{verbatim}
int binary_search(const std::vector<int>& a, int lo, int hi, int key) {
    if (lo > hi) return -1;                          // base: empty range
    int mid = lo + (hi - lo) / 2;                    // avoid overflow
    if (a[mid] == key)  return mid;
    if (a[mid] < key)   return binary_search(a, mid + 1, hi, key);
    else                return binary_search(a, lo, mid - 1, key);
}

// Caller:
int idx = binary_search(v, 0, v.size() - 1, 42);
\end{verbatim}
\end{examplebox}

\begin{warnbox}[\texttt{mid = (lo + hi) / 2} is an overflow bug]
If \texttt{lo + hi} exceeds \texttt{INT\_MAX}, the sum wraps negative
and \texttt{mid} lands somewhere random. The fix is
\texttt{mid = lo + (hi - lo) / 2}, which computes the same value
without ever adding two large indices. This bug sat undiscovered in
the JDK's \texttt{Arrays.binarySearch} for nearly a decade (Josh
Bloch, 2006). Make the safe form muscle memory.
\end{warnbox}

\subsection{Tracing binary search on a concrete input}

\begin{examplebox}[Search for 42 in {\normalfont (14, 26, 42, 59, 71, 88, 92)}]
\begin{verbatim}
bs(v, 0, 6, 42):   mid = 3,  v[3] = 59.  59 > 42  ->  search left
bs(v, 0, 2, 42):   mid = 1,  v[1] = 26.  26 < 42  ->  search right
bs(v, 2, 2, 42):   mid = 2,  v[2] = 42.  found. return 2
\end{verbatim}
Three comparisons to find the key in a 7-element list. A linear scan
would also take 3 here, but scale the list to $10^9$ elements: linear
is $\sim 10^9$ comparisons, binary is $\sim 30$. That's the
exponential-vs-logarithmic gap.
\end{examplebox}

\begin{notebox}[Why $O(\log n)$]
Each recursive call halves the remaining range. Starting from $n$
elements, you reach a range of size 1 after $\log_2 n$ halvings. 7
elements $\to$ $\log_2 7 \approx 3$ steps. $10^9$ elements $\to$ 30
steps. Binary search is the archetypal divide-and-conquer algorithm.
\end{notebox}

\subsection{Preconditions and edge cases}

\begin{warnbox}[The array MUST be sorted]
Binary search on an unsorted array silently returns wrong answers,
not an error. If the data isn't already sorted, either sort first
(which costs $O(n \log n)$ and only pays off when you will search
many times) or use linear search.
\end{warnbox}

\begin{warnbox}[The off-by-one on the recursive range]
The recursive calls use \texttt{mid + 1} and \texttt{mid - 1}, not
\texttt{mid}. Forgetting the $\pm 1$ leaves \texttt{mid} in the
remaining range forever, and when \texttt{a[mid]} keeps failing to
match, the recursion never terminates. This is exactly the ``make
the subproblem strictly smaller'' rule from \textsection 2.4.
\end{warnbox}

\begin{notebox}[\texttt{std::binary\_search} and \texttt{std::lower\_bound}]
The standard library already ships binary search, and the iterative
version it uses is somewhat faster in practice:
\begin{verbatim}
#include <algorithm>
bool found = std::binary_search(v.begin(), v.end(), key);

auto it = std::lower_bound(v.begin(), v.end(), key);
// it points to the first element >= key
\end{verbatim}
Use these in production code; write the recursion by hand to
understand it.
\end{notebox}

\subsection{Side-by-side: when recursion helps and when it hurts}

\begin{center}
\small
\begin{tabular}{lll}
& \textbf{Fibonacci (naive)} & \textbf{Binary search} \\
\hline
Recursive calls per step & 2 & 1 \\
Subproblem overlap & massive & none \\
Time complexity & $O(\phi^n)$ & $O(\log n)$ \\
Typical depth & $n$ & $\log n$ \\
Stack risk at $n = 10^6$ & overflow long before & safe (depth $\sim 20$) \\
\end{tabular}
\normalsize
\end{center}

\begin{ideabox}[Big picture]
Recursive definitions are free; efficient recursive algorithms are
not. When the recursion tree has \emph{branching plus overlap}, you
pay exponentially -- Fibonacci is the small-but-brutal example.
When each call strictly shrinks the problem and doesn't duplicate
work, recursion is often the clearest possible code -- binary search
is the archetype. A useful habit when writing a recursive function:
sketch the call tree for a small input. If the tree is a single
chain or a balanced split, the algorithm is fine. If you see the
same subproblem appearing on multiple branches, stop and think about
memoization or a loop before writing code.
\end{ideabox}

\section{2.6 Data Privacy}

This section interrupts the algorithms thread to cover professional
ethics -- specifically, how a working developer should think about
confidential data. The material is short on C++ but the habits it
teaches show up in every real job: which data you collect, where you
store it, who can see it, and how to reason about the tradeoffs when
those choices conflict.

\begin{defnbox}[Confidential data]
\textbf{Confidential data} is any data whose exposure could harm the
people or organization it describes -- names, home addresses, Social
Security numbers, health records, passwords, financial records,
written communications, location history. The baseline assumption
must be that unauthorized people will at some point try (or
accidentally manage) to see it.
\end{defnbox}

\begin{defnbox}[Data breach]
A \textbf{data breach} is any incident in which confidential data is
exposed to an unauthorized party -- by external attack, insider
misuse, accidental disclosure, physical theft, or lost hardware.
Breaches are rarely caused by a single dramatic failure; they're the
compounding effect of many small lapses.
\end{defnbox}

\subsection{``Alertness'' as a professional habit}

The textbook framing borrows a line from safety engineering: the worst
failures are the ones a little alertness would have prevented. For a
developer, ``alertness'' translates to a handful of reflexes:

\begin{itemize}[leftmargin=*]
    \item Before writing code that collects a piece of data, ask
          whether it's actually needed.
    \item Before sending a file, email, or dataset, ask who else sees
          it along the way.
    \item Before stepping away from a laptop with sensitive data open,
          lock the screen.
    \item Before committing, check for credentials, keys, or customer
          data in the diff.
\end{itemize}

None of these are technical. All of them are prevented by a quick
pause, and none of them take a second when they become habit.

\subsection{The Company XYZ scenario (and how it should have gone)}

The textbook scenario: a data analyst is hired to do a confidential
job-satisfaction survey. Data collected: name, home address,
supervisor, free-form criticism. Data stored: an unencrypted Excel
spreadsheet on a laptop left unlocked in a break room. Data leaked:
the spreadsheet, forwarded by an admin who saw it open. Fallout:
harassment, lawsuit, resignations.

\begin{examplebox}[Four separate failures, any one of which would have prevented it]
\begin{itemize}[leftmargin=*,noitemsep]
    \item \textbf{Collected too much.} Home address wasn't needed for
          ``satisfaction correlated with commute'' -- ZIP code or
          commute-time bucket would have worked.
    \item \textbf{Stored it identified.} Responses could have been
          anonymized at intake (swap name for a random ID, discard the
          mapping after collection closes).
    \item \textbf{Stored it in the clear.} The file should have been
          encrypted at rest and protected by a strong password.
    \item \textbf{Left the device unlocked.} A password-locked screen
          and a full-disk-encrypted laptop would have kept the casual
          snooper out.
\end{itemize}
Each of these is minutes of work. Together they turn a career-ending
breach into a non-event.
\end{examplebox}

\subsection{The two levers: collection and release}

\begin{defnbox}[Minimize collection risk]
Collect the \emph{least} data that answers the question. If a
summary suffices, collect the summary. If an anonymous identifier
suffices, don't collect the name. Data that was never collected
cannot leak.
\end{defnbox}

\begin{defnbox}[Minimize release risk]
Treat every piece of collected data as if it will eventually be
exposed. Encrypt at rest, lock access behind authentication, partition
storage so a single breach doesn't expose everything, and delete data
when it's no longer needed.
\end{defnbox}

\begin{notebox}[The eBay and Target breaches show the second principle at work -- partially]
eBay's 2014 breach exposed names, addresses, birth dates, and
encrypted passwords \emph{because financial data was stored
separately} and therefore did not leak with the rest. Target's 2013
breach did include credit-card numbers because the same segmentation
was not in place. Same kind of attack; different architectural
decisions about what data lives where; very different outcomes.
\end{notebox}

\subsection{The tradeoffs are real}

There is no zero-risk data practice. Every privacy protection costs
something -- inconvenience, reduced analytic power, slower development
velocity -- and the job is to \emph{consciously} balance cost against
risk.

\begin{center}
\small
\begin{tabular}{p{4.8cm} p{7.2cm}}
\textbf{Precaution} & \textbf{Tradeoff} \\
\hline
Drop identifying fields & Lose some analyses (commute distance,
per-supervisor complaint patterns). \\
Encrypt the file & Extra step every time you open it; password must
be remembered and rotated. \\
Collect fewer responses & Smaller sample, less statistical power. \\
Refuse to share a dataset & Slows collaboration; may block a
legitimate use. \\
Retain data longer & More value for longitudinal analysis, more
exposure if breached. \\
\end{tabular}
\normalsize
\end{center}

\begin{ideabox}[A four-question checklist]
Before collecting, storing, or sharing anything sensitive:
\begin{enumerate}[leftmargin=*,noitemsep]
    \item \textbf{Is this analysis worth the exposure risk if this
          data leaks?}
    \item \textbf{Can I avoid collecting (or can I anonymize) the
          sensitive fields?}
    \item \textbf{Is this data protected at rest} -- encryption,
          access control, locked device?
    \item \textbf{If someone else wants this data, do they
          \emph{need} it -- or would a summary suffice?}
\end{enumerate}
These four questions catch most breaches before they happen.
\end{ideabox}

\subsection{A short history of things that went wrong}

The scale of real breaches is hard to picture:

\begin{center}
\small
\begin{tabular}{p{2.2cm} p{9.8cm}}
\hline
\textbf{2013, Target} & 70M+ customer records including credit-card
data; stemmed from a compromised HVAC vendor login. \\
\textbf{2013, Snowden / NSA} & Thousands of classified documents
leaked by an insider with legitimate access. \\
\textbf{2014, eBay} & 148M accounts; passwords encrypted but email
and address in the clear. \\
\textbf{2014, Sony Pictures} & Terabytes including emails, salaries,
and HR data; a senior executive lost their job over quoted emails. \\
\textbf{2015, Anthem} & 80M records including SSNs and health info. \\
\textbf{2015, Ashley Madison} & 30M user identities of a service
whose entire value was discretion; several suicides attributed in
part to the exposure. \\
\textbf{2017, Equifax} & 147M Americans' SSNs and credit data; root
cause a known Apache Struts vulnerability left unpatched. \\
\textbf{2021, Facebook} & 533M user records including phone numbers
scraped via a contact-import misuse. \\
\hline
\end{tabular}
\normalsize
\end{center}

\begin{notebox}[The pattern in every row]
Large, well-resourced companies with security teams still got
breached. The common thread isn't ``they were incompetent'' -- it's
that attack surface is huge, humans make small mistakes, and modern
datasets are valuable enough that attackers are patient. Assuming
breach is the default working assumption; the question is how much
damage a breach would do, not whether one will ever happen.
\end{notebox}

\subsection{What this means for you on CS 300 and beyond}

Course projects rarely handle real PII, but the habits are worth
forming now:

\begin{itemize}[leftmargin=*,noitemsep]
    \item Don't commit credentials, API keys, or \texttt{.env} files
          to git. Add them to \texttt{.gitignore} before you run your
          first \texttt{git add}.
    \item Use environment variables or a secrets manager for
          API tokens -- never string literals in source.
    \item When a program reads user data, think about where it lives
          after the program exits: in memory? in a temp file? in a
          log?
    \item Treat student test data with the same care you would real
          data, so the habits hold when the data does matter.
\end{itemize}

\begin{ideabox}[Big picture]
Every algorithm in this course is a \emph{means}, and data is what
the algorithms operate on. Efficiency and correctness are necessary
but not sufficient -- correctness on the wrong data, efficiency at
the wrong scope, or a correct analysis leaked to the wrong audience
can all cause real harm. Professional responsibility isn't a separate
subject from engineering; it's the part of engineering that keeps
your correct code from becoming a liability.
\end{ideabox}

\section{2.7 Ethical Guidelines}

\textsection 2.6 covered the narrow question of keeping confidential
data confidential. This section widens the frame: once you have data
and are doing something with it, what obligations come with the work?
The textbook reaches for the American Statistical Association's
\textbf{Ethical Guidelines for Statistical Practice}, which is a good
reference point even though CS 300 is not a statistics course --
every engineer who ships a system that affects people is doing applied
statistics whether they call it that or not.

\begin{defnbox}[Ethics (working definition)]
\textbf{Ethics} are the principles that guide decisions about what is
good for individuals and society. Professional ethics are the
specific rules a field commits to so that its members can be trusted
to do their work responsibly -- analogous to the Hippocratic oath in
medicine or a building code in civil engineering.
\end{defnbox}

\begin{defnbox}[Data analytics]
The practice of examining data to draw conclusions, build models, or
inform decisions. Results are often used by people who did not
perform the analysis and cannot independently verify it, which is
why transparency and integrity matter so much.
\end{defnbox}

\subsection{The one-line summary}

The ASA opens its guidelines with a sentence worth memorizing:
\begin{quote}
Good statistical practice is fundamentally based on \textbf{transparent
assumptions}, \textbf{reproducible results}, and \textbf{valid
interpretations}.
\end{quote}

Those three phrases are the compressed form of everything that
follows. When a result looks impressive, ask: were the assumptions
stated? can someone else reproduce it? does the interpretation
actually follow from the numbers?

\subsection{The five categories}

The ASA guidelines organize obligations into five buckets. Roughly:
what you owe your work, your data, your audience, your subjects, and
your field.

\paragraph{1. Professional integrity and accountability.}
\begin{itemize}[leftmargin=*,noitemsep]
    \item Pick the sampling and analytic approach that actually
          answers the question, not the one that gives the answer
          you wanted.
    \item Identify and reduce sources of bias. Disclose what you
          can't eliminate.
    \item Disclose conflicts of interest -- financial, political, or
          otherwise. Readers get to discount your work accordingly.
\end{itemize}

\paragraph{2. Integrity of data and methods.}
\begin{itemize}[leftmargin=*,noitemsep]
    \item Document what you did. Not ``document in principle'' -- the
          preprocessing, the filters, the parameters, the cutoffs.
    \item Report results truthfully, including the uncertainty.
          ``Our model is 94\% accurate'' is meaningless without a
          confidence interval, a baseline, and the data it was
          measured on.
    \item Share data and code when you can, redacting or aggregating
          to preserve privacy when you can't.
\end{itemize}

\paragraph{3. Responsibilities to the public and/or client.}
\begin{itemize}[leftmargin=*,noitemsep]
    \item Explain the tradeoffs of the approach you chose, including
          what it cannot tell them.
    \item Make useful knowledge broadly available when it doesn't
          compromise privacy or legitimate confidentiality.
\end{itemize}

\paragraph{4. Responsibilities to research subjects.}
\begin{itemize}[leftmargin=*,noitemsep]
    \item Stay current on best practices for protecting the people
          whose data you hold.
    \item Protect privacy and confidentiality of subjects even when
          doing so costs you analytic power.
\end{itemize}

\paragraph{5. Allegations of misconduct.}
\begin{itemize}[leftmargin=*,noitemsep]
    \item Don't condone -- or cover for -- bad practice by
          colleagues.
    \item Distinguish honest error and genuine disagreement from
          misconduct. They look similar and are judged very
          differently.
    \item Don't retaliate against whistleblowers.
\end{itemize}

\subsection{Connecting the categories to CS 300 work}

``Statistical practice'' sounds distant from a sophomore algorithms
class. It shouldn't. A few concrete mappings:

\begin{center}
\small
\begin{tabular}{p{5cm} p{7cm}}
\textbf{Guideline} & \textbf{What it looks like in CS 300 / future work} \\
\hline
Transparent assumptions & Document your time complexity, your input
size assumptions, and what ``works'' means (average? worst case?
specific distribution?). \\
Reproducible results & Fixed random seeds in tests, deterministic
hash-table iteration, version-pinned dependencies. \\
Valid interpretation & ``Sort X finished 2$\times$ faster'' is a
claim about the specific input, machine, and compiler flags --
not a universal truth. \\
Disclose conflicts & If you generated your test data with the same
code you're benchmarking, say so. \\
Protect subjects & Don't commit student records, user IDs, or real
customer data to your course repo. \\
Don't retaliate, don't condone & If a teammate ships a function
that rounds metric values ``to look better,'' call it out in
review. \\
\end{tabular}
\normalsize
\end{center}

\subsection{The classic ways analyses go bad}

Unethical analysis often doesn't look like fraud -- it looks like a
slightly wrong choice at each of five steps that compound into a
misleading conclusion.

\begin{warnbox}[\emph{p}-hacking / selective reporting]
Running many analyses and reporting only the ones that came out
``interesting.'' The reported result looks significant; the
underlying process was essentially random. Pre-register the analysis
you intend to run before you look at the data, and report all runs,
not just the one that worked.
\end{warnbox}

\begin{warnbox}[Cherry-picked baselines and time windows]
``Model X improved accuracy 20\%'' -- compared to what, and over
which time window? Picking the weakest baseline or the favorable
date range can make an unchanged algorithm look transformative. The
mitigation is always to report the alternatives you considered and
why you chose this one.
\end{warnbox}

\begin{warnbox}[Disclosing conclusions, hiding methods]
A result that cannot be reproduced is not a result. If a paper,
blog post, or dashboard gives you a number without enough detail to
recompute it, treat the number as a claim, not evidence.
\end{warnbox}

\begin{warnbox}[Proxy variables for protected attributes]
A model that doesn't take race/gender/etc.\ as an input but uses ZIP
code, first name, or browsing history can reproduce discrimination
with a clean deniability story. The ethical question is not what
features you included but what outcomes the system produces across
groups.
\end{warnbox}

\subsection{A short CS 300 reproducibility checklist}

Even in course projects, form the habits now:

\begin{itemize}[leftmargin=*,noitemsep]
    \item \textbf{Seed} any randomness you use in benchmarks or tests.
    \item \textbf{Record} compiler, flags, CPU, and input size alongside
          every performance number.
    \item \textbf{Version} the dataset. ``The bids.csv we were given''
          is not enough -- keep a hash or a dated copy.
    \item \textbf{Rerun} the whole pipeline from scratch occasionally
          to make sure it still produces the reported results.
    \item \textbf{Commit} the analysis code next to the numbers it
          produced, not separately.
\end{itemize}

\begin{ideabox}[Big picture]
Ethical practice is not a separate module bolted onto ``the real
work.'' The three pillars -- transparent assumptions, reproducible
results, valid interpretations -- are the same hygiene that makes
engineering teams productive: code you can rerun, tests you can
trust, and claims that survive scrutiny. A benchmark nobody else
can reproduce is closer to marketing than to engineering; a model
whose behavior across subgroups nobody bothered to check is a bug
waiting to surface. Treat ethics as a quality bar you hold your own
work to, and the ASA guidelines stop feeling like an outside
constraint and start feeling like the definition of ``done.''
\end{ideabox}

\section{2.8 Heuristics}

\textsection 2.1 ended on a note: some problems (the NP-complete ones)
have no known efficient exact algorithm. What do you do when you still
have to ship? You use a \textbf{heuristic} -- an algorithm that
deliberately trades optimality for speed or simplicity. Heuristics
are everywhere in production code: route planners, ranking systems,
allocators, schedulers, compilers, game AI. When the exact answer is
too expensive to compute, a provably-pretty-good answer in a
reasonable amount of time is the whole game.

\begin{defnbox}[Heuristic]
A \textbf{heuristic} is a technique that willingly accepts a
non-optimal or approximate solution in order to improve execution
speed, simplicity, or memory use. A heuristic is not a
``wrong'' algorithm -- it's a conscious tradeoff. The tradeoff is the
point.
\end{defnbox}

\begin{defnbox}[Heuristic algorithm]
An algorithm whose design is driven by a heuristic. It typically runs
much faster than an optimal algorithm on the same problem, and
typically returns a solution that is \emph{close to} optimal most of
the time -- with no guarantee, or a weaker guarantee, of hitting the
best possible answer.
\end{defnbox}

\subsection{The knapsack problem: the canonical example}

\begin{defnbox}[0-1 knapsack problem]
Given a knapsack with maximum weight capacity $W$ and a list of
items, each with a weight $w_i$ and a value $v_i$, choose a subset of
items (each included at most once -- hence \textbf{0-1}) whose total
weight is $\le W$ and whose total value is as large as possible.
\end{defnbox}

\begin{notebox}[Why this is hard]
With $n$ items, there are $2^n$ possible subsets. Checking all of
them is exact but exponential. 30 items is already $2^{30}
\approx 10^9$ combinations; 50 items is $10^{15}$. This is the
textbook example of ``correct algorithm exists, but good luck
running it.'' 0-1 knapsack is NP-hard; no polynomial-time exact
algorithm is known.
\end{notebox}

\subsection{A greedy heuristic for knapsack}

\begin{examplebox}[Greedy by value]
\begin{verbatim}
Knapsack01_greedy(items, capacity):
    Sort items in descending order of value
    remaining = capacity
    chosen = empty list
    for each item in sorted order:
        if item.weight <= remaining:
            add item to chosen
            remaining -= item.weight
    return chosen
\end{verbatim}
Take the highest-value item that still fits; keep going until nothing
else fits. One pass, plus the sort: $O(n \log n)$ total.
\end{examplebox}

\begin{warnbox}[This is fast, and sometimes wrong]
\textbf{Capacity 20, items:}
\begin{center}
\begin{tabular}{lll}
Item & Weight & Value \\
\hline
A & 15 & \$95 \\
B & 12 & \$60 \\
C & \phantom{0}8 & \$42 \\
\end{tabular}
\end{center}
Greedy picks A (\$95, remaining capacity 5) and then nothing else
fits. Total: \$95. The optimal answer picks B + C: weight 20, value
\$102. The heuristic missed the best answer by choosing the
tempting-looking big-ticket item early.
\end{warnbox}

\begin{ideabox}[Smarter greedy: value per pound]
Sorting by \emph{value/weight ratio} instead of raw value usually
does better, because it packs the knapsack more densely. It still
isn't optimal on every instance -- the same pathological example
can be constructed against any simple greedy rule -- but on random
instances it's typically within a few percent of optimal. Choosing
the right heuristic measure is half the design work.
\end{ideabox}

\begin{notebox}[The other two options: DP and branch-and-bound]
If all weights are integers and $W$ is not too large, 0-1 knapsack
can be solved \emph{exactly} in $O(n \cdot W)$ time with dynamic
programming. For large $W$ or continuous weights, branch-and-bound
prunes the exponential search tree but retains the correct answer.
The full story of ``when to use heuristic vs.\ DP vs.\
branch-and-bound'' is beyond this section, but knowing that all
three exist helps you reach for the right tool.
\end{notebox}

\begin{notebox}[``DP solves knapsack'' is true -- but only \emph{pseudopolynomially}]
The line above says 0-1 knapsack can be solved in $O(n \cdot W)$
time with DP, which sounds polynomial. It's not -- not in the sense
algorithm theorists mean.

\textbf{Strongly polynomial} time is bounded by a constant-degree
polynomial in input size measured \emph{in words}. Knapsack's input
is $n$ items plus the integer $W$ -- so input size is $O(n)$ words.
But $W$ itself is one of those words: it can be as large as $2^w$
where $w$ is the word size. So $O(n \cdot W)$ can be $O(n \cdot 2^w)$
-- exponential in input \emph{bits}.

This is called \textbf{pseudopolynomial}: bounded by a polynomial in
input size \emph{and} input integer values. Pseudopolynomial
algorithms are polynomial only when the integers are polynomially
bounded in input size ($n^{O(1)}$). For knapsack with $W = 2^{60}$,
$O(n \cdot W)$ is hopeless.

So 0-1 knapsack remains NP-hard -- the DP doesn't actually break
the hardness. The DP is useful when $W$ is small (a few thousand,
say), but fundamentally the difficulty sits in the \emph{values} of
the integers, not their count. Subset Sum has the same character
(same DP, $O(n \cdot T)$, pseudopolynomial in target $T$). Chapter~3
returns to the strongly-polynomial-vs-pseudopolynomial-vs-exponential
hierarchy when it formalizes Big-O.
\end{notebox}

\subsection{Self-adjusting heuristics}

A different flavor of heuristic: instead of approximating an answer,
the algorithm \emph{reorganizes its data structure over time} based
on how it's being used, so common operations get faster.

\begin{defnbox}[Self-adjusting heuristic]
A heuristic that modifies a data structure in response to observed
access patterns. The structure starts in some default arrangement and
evolves to match the workload. Splay trees, AVL rotations, LRU caches,
move-to-front lists, and adaptive B-trees are all examples.
\end{defnbox}

\subsection{Worked example: move-to-front on a linear search}

\begin{examplebox}[Move-to-front]
Start with a list: \texttt{[17, 5, 88, 3, 42, 91, 76, 64]}.

Search for 42: linear scan takes 5 comparisons. After the search,
\emph{move 42 to the front}.

List is now: \texttt{[42, 17, 5, 88, 3, 91, 76, 64]}.

Search for 42 again: 1 comparison.

Search for 64: 8 comparisons, then move 64 to the front.

List: \texttt{[64, 42, 17, 5, 88, 3, 91, 76]}.

Search for 42: 2 comparisons.
\end{examplebox}

\begin{ideabox}[Why it works]
Real-world access patterns are rarely uniform. A small number of
items are hot, most are cold. Move-to-front drifts the hot items
toward the front of the list, where they're cheap to find. On a
skewed workload -- think the handful of popular queries in a search
engine, or the same test scores being looked up hundreds of times --
this can make an $O(n)$ linear search effectively $O(1)$ amortized.

On a \emph{uniform} workload, move-to-front still works, it just
doesn't help much. That's the right kind of tradeoff: it never
makes access slower (per operation the list still takes $O(n)$
worst case) and it makes the common case much faster when the
common case exists.
\end{ideabox}

\begin{warnbox}[Move-to-front is not a substitute for the right data structure]
If the list is big and the workload isn't heavily skewed, a hash
table (\textsection later in course) is $O(1)$ average regardless
of the access pattern and you should use one. Move-to-front is a
nice $O(1)$-amortized tweak when a linked list is already the right
structure for other reasons.
\end{warnbox}

\subsection{When to reach for a heuristic (and when not to)}

\begin{center}
\small
\begin{tabular}{p{5.2cm} p{7cm}}
\textbf{Reach for a heuristic when} & \textbf{Use the exact algorithm when} \\
\hline
The problem is NP-hard and inputs are large. & The problem has a
polynomial-time exact algorithm (most of this course). \\
``Pretty good, fast'' is more valuable than ``optimal, slow.'' &
Optimality is required by the spec (safety, financial, legal). \\
You can measure how close the heuristic gets to optimal and it's
good enough. & You can't measure closeness and a bad answer has
real consequences. \\
The algorithm will run on streaming / online data where you can't
see the whole input. & The input is fully available and small
enough to process exactly. \\
\end{tabular}
\normalsize
\end{center}

\begin{ideabox}[Big picture]
Heuristics are not a failure of computer science -- they are one of
its honest responses to the real world. Some problems are too big to
solve optimally; others arrive as a stream where you have to commit
before seeing everything; many data structures get dramatically
faster when allowed to reshape themselves around actual usage. The
habit of mind to build here is: \emph{ask what ``good enough'' means
for this problem, pick a heuristic whose failure modes you
understand, and measure how close to optimal it gets in practice.}
That framing will serve you long after the specific greedy-knapsack
or move-to-front trick stops mattering.
\end{ideabox}

\section{2.9 Greedy Algorithms}

\textsection 2.8 introduced heuristics -- the idea of accepting
non-optimal answers for speed. A \textbf{greedy algorithm} is the most
common style of heuristic: at each step, commit to the choice that
looks best \emph{right now}, without looking ahead. Sometimes this
lands you at the globally optimal answer (change-making,
fractional-knapsack, activity selection, shortest paths). Sometimes
it doesn't (0-1 knapsack from \textsection 2.8). Knowing which
problems are in the first bucket and which aren't is the real skill.

\begin{defnbox}[Greedy algorithm]
A \textbf{greedy algorithm} solves a problem by repeatedly making the
\emph{locally optimal} choice -- the option that looks best at this
step -- and never reconsidering. It does not backtrack, cache
alternatives, or look ahead. The algorithm terminates when no further
choice is possible or the goal has been reached.
\end{defnbox}

\begin{defnbox}[Locally optimal vs. globally optimal]
A \textbf{locally optimal} choice is the best move at a single step
considered in isolation. A \textbf{globally optimal} solution is the
best overall outcome across all steps. Greedy algorithms make locally
optimal choices and \emph{hope} those add up to a global optimum.
Whether they do depends on the problem.
\end{defnbox}

\subsection{Example 1: making change (and why US coins are special)}

\begin{examplebox}[\texttt{MakeChange} with US denominations]
\begin{verbatim}
MakeChange(amount):
    while amount >= 25: take a quarter; amount -= 25
    while amount >= 10: take a dime;    amount -= 10
    while amount >= 5:  take a nickel;  amount -= 5
    while amount >= 1:  take a penny;   amount -= 1
\end{verbatim}
For 91 cents: 3 quarters (75), 1 dime (85), 1 nickel (90), 1 penny
(91). Total: 6 coins. That also happens to be the fewest coins
possible.
\end{examplebox}

\begin{warnbox}[Greedy change-making is NOT always optimal]
It works for US coins -- 1, 5, 10, 25 -- because the denomination set
has a property called \emph{matroid-compatibility}; for every amount,
grabbing the largest coin that fits is part of some optimal solution.

Change the denominations and greedy breaks. Example: coins
\{1, 3, 4\}, target 6. Greedy takes 4, then 1, then 1 -- three coins.
The optimal is 3 + 3 -- two coins. Same algorithm, different coin
set, different answer.

The exact-optimal solution for arbitrary denominations is dynamic
programming, $O(n \cdot A)$ in the amount $A$.
\end{warnbox}

\subsection{Example 2: fractional knapsack -- greedy is optimal}

\begin{defnbox}[Fractional knapsack]
The knapsack problem, but you may take any fraction $f \in [0, 1]$
of any item. Its weight contribution is $f \cdot w_i$ and its value
contribution is $f \cdot v_i$.
\end{defnbox}

\begin{examplebox}[Greedy by value-per-pound]
\begin{verbatim}
FractionalKnapsack(items, capacity):
    Sort items descending by (value / weight)
    remaining = capacity
    for each item in sorted order:
        if item.weight <= remaining:
            take all of item
            remaining -= item.weight
        else:
            take fraction = remaining / item.weight of item
            remaining = 0
            break
\end{verbatim}
This is \textbf{provably optimal}. Intuition: at every pound of
knapsack space, you prefer the item with the highest
dollars-per-pound. Allowing fractions means you can always
\emph{fill} the knapsack, so ``best value density packed first'' wins.
0-1 knapsack lacks this property because you're forced into an
all-or-nothing choice at the last item.
\end{examplebox}

\begin{ideabox}[The difference between 0-1 and fractional]
Same items, same capacity, same greedy rule -- but on 0-1 the greedy
can leave optimal value on the table (see \textsection 2.8), while on
fractional it can't. The lesson: a greedy rule's correctness depends
on the problem structure, not just the rule. Always ask ``is this
problem one where greedy is known to be optimal?'' before shipping
code.
\end{ideabox}

\subsection{Example 3: activity selection -- and why ``earliest finish'' wins}

\begin{defnbox}[Activity selection problem]
Given activities each with a start time and a finish time, pick the
largest subset of activities such that no two chosen activities
overlap. Equivalent to scheduling the maximum number of compatible
meetings on a single room.
\end{defnbox}

\begin{examplebox}[Greedy by earliest finish]
\begin{verbatim}
ActivitySelection(A):
    Sort A ascending by finish time
    chosen = [A[0]]
    current_end = A[0].finish
    for i from 1 to n-1:
        if A[i].start >= current_end:
            chosen.append(A[i])
            current_end = A[i].finish
    return chosen
\end{verbatim}
\end{examplebox}

\begin{notebox}[Why earliest finish, not shortest or earliest start]
Picking the shortest activity fails: a 10-minute activity that spans
lunch blocks both morning and afternoon longer ones. Picking the
earliest start fails for the same reason -- a 9am-to-5pm activity
beats out everything. \textbf{Earliest finish} leaves the maximum
remaining room in the day for further choices, which is exactly what
greedy needs to stay on track.
\end{notebox}

\begin{examplebox}[A vacation day worked by hand]
Activities: museum (9--11), hike (10--12), boat (11--13), hang-glide
(14--15), fireworks (19--20), nightclub (18--23).

Sorted by finish: museum, hike, boat, hang-glide, fireworks,
nightclub.

\begin{itemize}[leftmargin=*,noitemsep]
    \item Choose museum (finishes 11).
    \item Hike starts 10 < 11, skip.
    \item Boat starts 11 $\ge$ 11, choose (finishes 13).
    \item Hang-glide starts 14 $\ge$ 13, choose (finishes 15).
    \item Fireworks starts 19 $\ge$ 15, choose (finishes 20).
    \item Nightclub starts 18 < 20, skip.
\end{itemize}
Chosen: museum, boat, hang-glide, fireworks. That's 4 activities,
which is the maximum possible.
\end{examplebox}

\begin{examplebox}[Why earliest-finish is optimal: a sketch]
Activity selection is one of the few greedy algorithms with a
clean, short correctness proof. The 2-step template from above made
concrete:

\textbf{Greedy-choice property.} Some optimal solution includes the
activity with the earliest finish time. \emph{Why:} let
$\mathcal{O}$ be any optimal schedule, and let $a^*$ be the
earliest-finish activity overall (call its finish time $f^*$). If
$a^* \in \mathcal{O}$, done. Otherwise let $a' \in \mathcal{O}$ be
the activity in $\mathcal{O}$ with the earliest finish time. Since
$a^*$ has the globally earliest finish, $f^* \le f(a')$. Swap $a'$
out and $a^*$ in: the rest of $\mathcal{O}$ all started \emph{after}
$f(a') \ge f^*$, so they're still compatible with $a^*$. The new
schedule has the same number of activities, includes $a^*$, and is
also optimal. This is an \textbf{exchange argument} -- the standard
shape of a greedy-choice proof.

\textbf{Optimal substructure.} After picking $a^*$, the remaining
problem is: pick the most activities from those that start at or
after $f^*$. That subproblem has the same shape as the original
(still an activity-selection problem, on a smaller set). Combining
the greedy choice $a^*$ with an optimal solution to the remaining
subproblem gives an optimal solution to the whole.

\textbf{Putting them together.} Greedy-choice says picking $a^*$
loses nothing. Optimal substructure says we can recurse on the
remaining problem with the same algorithm. By induction on input
size, the algorithm produces an optimal schedule.

This is the entire shape of every greedy-correctness proof you
will write: an exchange argument for greedy-choice, plus a
same-shape argument for optimal substructure. Compare to the longer
DP proofs in \textsection 2.10 -- greedy is simpler when it
applies, and that simplicity is exactly why we look for
greedy-friendly structure first.
\end{examplebox}

\subsection{When greedy actually works}

There is a pattern behind these examples. Greedy is provably optimal
on problems with one of two underlying structures:

\begin{itemize}[leftmargin=*]
    \item \textbf{Matroid structure.} Fancy name, simple idea: if the
          locally optimal choice is always part of \emph{some}
          globally optimal solution, and ``independent subsets'' have
          an exchange property. Change-making with US coins,
          minimum-spanning-tree (Kruskal, Prim, later in course), and
          various scheduling problems have it.
    \item \textbf{Greedy-choice property $+$ optimal substructure.}
          Same two ingredients that justify dynamic programming,
          minus the need to enumerate every subproblem. Activity
          selection fits here.
\end{itemize}

\begin{ideabox}[Two moves of a correctness argument]
To argue a greedy algorithm is optimal you usually show:
\begin{enumerate}[leftmargin=*,noitemsep]
    \item \textbf{Greedy-choice property.} There exists an optimal
          solution that begins with the greedy choice. (``Swapping in
          the greedy item doesn't hurt.'')
    \item \textbf{Optimal substructure.} After making the greedy
          choice, the remaining problem has the same shape, and its
          optimal solution can be extended to a global optimum.
\end{enumerate}
Neither condition is obvious to verify; proving greedy correctness
is genuinely harder than writing the algorithm.
\end{ideabox}

\subsection{Greedy vs. other strategies}

\begin{center}
\small
\begin{tabular}{p{3.4cm} p{4cm} p{4.6cm}}
\textbf{Strategy} & \textbf{How it decides} & \textbf{Cost / guarantee} \\
\hline
Greedy & Best local choice, no reconsideration & Fast; optimal only
on certain problems \\
Dynamic programming & Remember optimal answers to subproblems, build
up & $O(n \cdot \text{state size})$; optimal when subproblems
overlap \\
Divide and conquer & Split, recurse, combine & Optimal; works when
subproblems are independent \\
Backtracking & Try choices, undo bad ones & Exponential worst case
but often prunable \\
Branch \& bound & Backtracking + bound pruning & Exact; can be
practical for small NP-hard inputs \\
\end{tabular}
\normalsize
\end{center}

\begin{warnbox}[The trap: assuming greedy works]
The most common mistake is to \emph{assume} the greedy algorithm
you wrote is optimal because it produced the right answer on the
example you tested. Try adversarial inputs. Read whether the
textbook problem is known to admit a greedy-optimal solution. When
in doubt, dynamic programming is usually within reach and
provably optimal.
\end{warnbox}

\begin{ideabox}[Big picture]
Greedy is the first algorithm-design strategy worth internalizing
because it is often the right answer and always worth trying first.
Write down the most natural ``best-at-each-step'' rule, test it
against an optimal solution on a few adversarial inputs, and either
(a) ship it with a proof sketch of why it's optimal, or (b) observe
it fail and fall back to DP or exact search. The three problems in
this section -- change, fractional knapsack, activity selection --
are worth keeping in your head as the canonical ``greedy wins''
references, and 0-1 knapsack from \textsection 2.8 as the canonical
``greedy loses'' reference.
\end{ideabox}

\section{2.10 Dynamic Programming}

Greedy (\textsection 2.9) makes one choice at each step and never
looks back. Pure recursion (\textsection 2.5) recomputes the same
subproblems over and over. \textbf{Dynamic programming (DP)} splits
the difference: it solves every subproblem \emph{once}, stores the
answer, and assembles the final answer from those stored pieces.

When the problem has overlapping subproblems and optimal
substructure, DP turns exponential-time recurrences into polynomial
algorithms -- usually a dramatic improvement.

\begin{defnbox}[Dynamic programming]
A problem-solving technique that splits a problem into smaller
subproblems, \textbf{computes each subproblem's answer exactly once},
stores it in a table, and uses the stored answers to build up the
final answer. Two common forms: \textbf{memoization} (top-down
recursion with a cache) and \textbf{tabulation} (bottom-up iteration
filling a table).
\end{defnbox}

\begin{defnbox}[Overlapping subproblems]
A recurrence has \textbf{overlapping subproblems} when the same
subproblem is computed multiple times in the naive recursion tree.
DP's payoff is skipping those repeated computations.
\end{defnbox}

\begin{defnbox}[Optimal substructure]
A problem has \textbf{optimal substructure} when an optimal solution
to the whole can be assembled from optimal solutions to its
subproblems. Without this, storing subproblem answers doesn't help.
\end{defnbox}

\subsection{Example 1: Fibonacci, fixed}

\textsection 2.5 showed naive recursive \texttt{fib} blowing up
exponentially because \texttt{fib(n-1)} and \texttt{fib(n-2)}
recompute overlapping work. DP fixes this two ways.

\paragraph{Top-down memoization.} Keep a cache, check it before
recursing:

\begin{examplebox}[Memoized recursion]
\begin{verbatim}
#include <vector>

std::vector<long long> memo;  // sized to n+1, filled with -1

long long fib(int n) {
    if (n <= 1) return n;
    if (memo[n] != -1) return memo[n];
    return memo[n] = fib(n - 1) + fib(n - 2);
}
\end{verbatim}
Each \texttt{fib(k)} gets computed once. $O(n)$ time, $O(n)$ memory
on the cache plus $O(n)$ on the call stack.
\end{examplebox}

\paragraph{Bottom-up tabulation.} Build the answers iteratively from
the base cases up:

\begin{examplebox}[Tabulated iteration]
\begin{verbatim}
long long fib(int n) {
    if (n <= 1) return n;
    long long prev = 0, cur = 1;
    for (int i = 2; i <= n; ++i) {
        long long next = prev + cur;
        prev = cur;
        cur = next;
    }
    return cur;
}
\end{verbatim}
$O(n)$ time, $O(1)$ memory, no recursion. This is usually the
preferred form when the dependency structure is simple.
\end{examplebox}

\begin{ideabox}[Top-down vs. bottom-up]
Top-down memoization is easier to write when a recurrence is already
in your head -- you write the natural recursion, add three lines of
caching, and stop. Bottom-up is usually tighter, avoids recursion
depth limits, and often lets you throw away old table entries you no
longer need (\textbf{rolling-array} trick -- the Fibonacci version
uses only two variables instead of an $n$-size array).
\end{ideabox}

\subsection{Example 2: longest common substring (LCS)}

\begin{notebox}[``LCS'' usually means subsequence, not substring]
\textbf{Heads up on naming.} The acronym \textbf{LCS} in most
textbooks (CLRS Ch.~15.4, MIT 6.006 Lecture 16, every DP reference
you'll Google) stands for \textbf{longest common \emph{subsequence}}
-- a non-contiguous match where characters from $s$ and $t$ must
appear in the same order but can have gaps in between. So
\texttt{ACE} is a common subsequence of \texttt{ABCDE} and
\texttt{XAYCEZ}.

What \emph{this} section covers is the longest common
\textbf{substring} (contiguous) variant -- a different problem with
a different recurrence:

\begin{center}
\small
\begin{tabular}{lll}
& \textbf{Substring (this section)} & \textbf{Subsequence (standard ``LCS'')} \\
\hline
Match must be          & contiguous          & in-order, gaps allowed \\
$L[i][j]$ if mismatch  & $0$ (reset)         & $\max(L[i{-}1][j],\ L[i][j{-}1])$ \\
Answer location        & max anywhere in table & $L[m][n]$ corner \\
\end{tabular}
\normalsize
\end{center}

When you see ``LCS'' outside this chapter -- in a textbook,
interview question, or Stack Overflow answer -- assume
\emph{subsequence} unless context says otherwise. Both variants are
useful; just don't confuse them.
\end{notebox}

A classic DP where the benefit is overwhelming. Naive recursion is
exponential; DP is $O(n \cdot m)$ in the two string lengths. This
algorithm underlies real-world tools: DNA sequence alignment, git's
rename detection, plagiarism detection.

\begin{defnbox}[Longest common substring problem]
Given two strings \texttt{s} and \texttt{t}, find the longest
\emph{contiguous} substring that appears in both. (Distinct from
longest common \emph{subsequence}, where characters don't have to be
adjacent.)
\end{defnbox}

\paragraph{The recurrence.} Let $L[i][j]$ be the length of the
longest common substring that \emph{ends at} \texttt{s[i]} and
\texttt{t[j]}.

\[
L[i][j] = \begin{cases}
0 & \text{if } s[i] \ne t[j] \\
1 + L[i-1][j-1] & \text{if } s[i] = t[j]
\end{cases}
\]

The answer is the maximum value anywhere in the table. Its position
tells you where the substring ends in \texttt{s}, which together with
the length gives you the substring itself.

\begin{examplebox}[LCS algorithm]
\begin{verbatim}
#include <string>
#include <vector>

std::string lcs(const std::string& s, const std::string& t) {
    int n = s.size(), m = t.size();
    std::vector<std::vector<int>> L(n, std::vector<int>(m, 0));

    int best_len = 0, best_row = 0;
    for (int i = 0; i < n; ++i) {
        for (int j = 0; j < m; ++j) {
            if (s[i] == t[j]) {
                L[i][j] = (i > 0 && j > 0) ? L[i-1][j-1] + 1 : 1;
                if (L[i][j] > best_len) {
                    best_len = L[i][j];
                    best_row = i;
                }
            }
        }
    }
    int start = best_row - best_len + 1;
    return s.substr(start, best_len);
}
\end{verbatim}
$O(n \cdot m)$ time, $O(n \cdot m)$ memory. For typical inputs
(dozens to thousands of characters) this is fine.
\end{examplebox}

\paragraph{Worked trace: \texttt{"Look"} vs \texttt{"Books"}.}
Rows indexed by \texttt{Look}, columns by \texttt{Books}.

\begin{center}
\small
\begin{tabular}{c|ccccc}
 & B & o & o & k & s \\
\hline
L & 0 & 0 & 0 & 0 & 0 \\
o & 0 & 1 & 1 & 0 & 0 \\
o & 0 & 1 & 2 & 0 & 0 \\
k & 0 & 0 & 0 & 3 & 0 \\
\end{tabular}
\normalsize
\end{center}

The maximum is 3, in row 3 (the \texttt{k}) column 3. The substring
ends at \texttt{s[3] = 'k'} and is 3 long: \texttt{"ook"}.

\subsection{Space optimization: the rolling-array trick}

The LCS recurrence only looks one row back. So the whole table is
wasteful -- you only need the previous row and the current row.

\begin{examplebox}[Two-row LCS]
\begin{verbatim}
std::vector<int> prev(m, 0), cur(m, 0);
int best_len = 0, best_row = 0;
for (int i = 0; i < n; ++i) {
    for (int j = 0; j < m; ++j) {
        if (s[i] == t[j])
            cur[j] = (j > 0 ? prev[j-1] : 0) + 1;
        else
            cur[j] = 0;
        if (cur[j] > best_len) { best_len = cur[j]; best_row = i; }
    }
    std::swap(prev, cur);
    std::fill(cur.begin(), cur.end(), 0);
}
\end{verbatim}
Same $O(n \cdot m)$ time. Memory drops from $O(n \cdot m)$ to
$O(m)$. On the billion-character DNA strings the textbook mentions,
this distinction is ``runs in memory'' vs.\ ``doesn't fit in RAM.''
\end{examplebox}

\subsection{Recognizing DP problems}

Not every recursion wants memoization. A few signals that DP is the
right tool:

\begin{itemize}[leftmargin=*]
    \item The naive recurrence recomputes the same subproblem on
          multiple branches (Fibonacci, edit distance, LCS, coin
          change, knapsack).
    \item The problem has a clear \emph{state} -- a small tuple of
          integers (index into a string, remaining capacity,
          whether you're allowed to take an item, etc.) -- and the
          answer depends only on that state.
    \item An optimal answer to the whole can be built by combining
          optimal answers to smaller states.
\end{itemize}

If the recursion tree doesn't overlap (binary search, merge sort),
plain recursion is fine and DP adds overhead for no benefit. If the
problem lacks optimal substructure (many game-theoretic or
greedy-fails problems), DP doesn't apply and you need a different
technique.

\begin{warnbox}[DP is not automatic -- designing the state is the hard part]
``Apply dynamic programming'' sounds like a step, but the real work
is identifying \emph{what subproblem to solve}. Most DP bugs come
from a state definition that either misses necessary information
(answer depends on something you didn't track) or includes
unnecessary information (table blows up exponentially). Spend the
thinking up front defining $L[i][j] = $ \emph{``the
\underline{\ \ \ \ \ \ \ \ \ \ }''} precisely; the code is almost
mechanical after that.
\end{warnbox}

\subsection{DP vs. greedy vs. divide-and-conquer}

\begin{center}
\small
\begin{tabular}{p{3.0cm} p{4.5cm} p{4.5cm}}
& \textbf{Greedy} & \textbf{DP} \\
\hline
Decision style & One choice per step, no backtrack & Considers all
choices for each subproblem \\
Subproblems overlap & (doesn't matter) & Required for DP to help \\
Optimal substructure needed & yes & yes \\
Guarantees & Optimal only for some problems & Always optimal when
it applies \\
Typical cost & $O(n)$ or $O(n \log n)$ & Product of state
dimensions \\
\end{tabular}
\normalsize
\end{center}

\begin{notebox}[Divide and conquer is the third cousin]
Divide-and-conquer (merge sort, binary search, Karatsuba) also
breaks a problem into subproblems, but the subproblems are
\emph{independent} -- no need to cache, no overlap. When the
recurrence happens to produce \emph{disjoint} subproblems, D\&C is
the natural form; when it produces \emph{overlapping} ones, DP
earns its keep.
\end{notebox}

\begin{ideabox}[Big picture]
DP is the natural upgrade from recursion when the recurrence tree
starts exploding. Every DP solution answers one question: ``what
state do I need to track so that the optimal answer depends only on
that state and the answers to smaller states?'' Once you've written
that down, the rest is filling a table. The two examples in this
section -- Fibonacci (one-dimensional state) and LCS
(two-dimensional state) -- are worth keeping in your head as the
shape of every DP solution you'll meet later. Coin change,
knapsack, edit distance, and most string-alignment problems follow
the same template.
\end{ideabox}

\section*{What this connects to}

\begin{notebox}[Forward links to later chapters]
\begin{itemize}
    \item \textbf{Chapter 3 (Big-O).} The informal ``na\"ive Fibonacci is
    exponential, memoized is linear'' claim becomes the precise
    statement $O(2^n)$ vs.\ $O(n)$. Same story for every algorithmic
    claim you made in this chapter -- ch.~3 supplies the notation.
    \item \textbf{Sorting in ch.~3.} Insertion and selection sort are
    incremental; merge sort and quicksort are divide-and-conquer
    (recursive); the runtime analyses reuse the recursion-tree reasoning
    from 2.5 verbatim.
    \item \textbf{Binary search} (2.5) will reappear as the inner loop
    of BST search (ch.~6) and as the \texttt{std::lower\_bound}
    primitive used everywhere.
    \item \textbf{Dynamic programming} (2.10) resurfaces in several
    later-course problems: longest-path-in-DAG, knapsack, edit
    distance, string alignment -- all use the same ``define the state,
    write the recurrence, fill the table'' workflow.
    \item \textbf{Greedy} (2.9) is the engine behind Dijkstra's algorithm
    and Kruskal's MST (optional ch.~10), and behind Huffman coding in
    compression. The greedy-choice property test you practiced here is
    exactly what justifies those algorithms' correctness.
    \item \textbf{Ethics / privacy} (2.6-2.7) isn't a detour -- it's the
    lens through which the course will ask you to evaluate every project.
    Expect discussion-board prompts on this all semester.
\end{itemize}
\end{notebox}

\textbf{Companion materials.} Compact notes:
\texttt{chapters/ch\_2/notes.tex}. Practice prompts (problem $\to$
pseudocode $\to$ C++ workflow): \texttt{chapters/ch\_2/practice.md}.

\end{document}
